%%
%% This is file `sample-sigconf-authordraft.tex',
%% generated with the docstrip utility.
%%
%% The original source files were:
%%
%% samples.dtx  (with options: `all,proceedings,bibtex,authordraft')
%% 
%% IMPORTANT NOTICE:
%% 
%% For the copyright see the source file.
%% 
%% Any modified versions of this file must be renamed
%% with new filenames distinct from sample-sigconf-authordraft.tex.
%% 
%% For distribution of the original source see the terms
%% for copying and modification in the file samples.dtx.
%% 
%% This generated file may be distributed as long as the
%% original source files, as listed above, are part of the
%% same distribution. (The sources need not necessarily be
%% in the same archive or directory.)
%%
%%
%% Commands for TeXCount
%TC:macro \cite [option:text,text]
%TC:macro \citep [option:text,text]
%TC:macro \citet [option:text,text]
%TC:envir table 0 1
%TC:envir table* 0 1
%TC:envir tabular [ignore] word
%TC:envir displaymath 0 word
%TC:envir math 0 word
%TC:envir comment 0 0
%%
%% The first command in your LaTeX source must be the \documentclass
%% command.
%%
%% For submission and review of your manuscript please change the
%% command to \documentclass[manuscript, screen, review]{acmart}.
%%
%% When submitting camera ready or to TAPS, please change the command
%% to \documentclass[sigconf]{acmart} or whichever template is required
%% for your publication.
%%
%%
\documentclass[sigchi,nonacm,manuscript]{acmart}

% For code listings used throughout the paper
\usepackage{listings}
\lstset{basicstyle=\ttfamily\small, breaklines=true, columns=fullflexible}

% For multirow tables
\usepackage{multirow}

% For subfigures (side-by-side images)
\usepackage{subcaption}
%%
%% \BibTeX command to typeset BibTeX logo in the docs
\AtBeginDocument{%
  \providecommand\BibTeX{{%
    Bib\TeX}}}

%% Rights management information.  This information is sent to you
%% when you complete the rights form.  These commands have SAMPLE
%% values in them; it is your responsibility as an author to replace
%% the commands and values with those provided to you when you
%% complete the rights form.
\setcopyright{acmlicensed}
\copyrightyear{2025}
\acmYear{2025}
\acmDOI{XXXXXXX.XXXXXXX}
%% These commands are for a PROCEEDINGS abstract or paper.
\acmConference[Conference acronym 'XX]{Make sure to enter the correct
  conference title from your rights confirmation email}{June 03--05,
  2018}{Woodstock, NY}
%%
%%  Uncomment \acmBooktitle if the title of the proceedings is different
%%  from ``Proceedings of ...''!
%%
%%\acmBooktitle{Woodstock '18: ACM Symposium on Neural Gaze Detection,
%%  June 03--05, 2018, Woodstock, NY}
\acmISBN{978-1-4503-XXXX-X/2018/06}


%%
%% Submission ID.
%% Use this when submitting an article to a sponsored event. You'll
%% receive a unique submission ID from the organizers
%% of the event, and this ID should be used as the parameter to this command.
%%\acmSubmissionID{123-A56-BU3}

%%
%% For managing citations, it is recommended to use bibliography
%% files in BibTeX format.
%%
%% You can then either use BibTeX with the ACM-Reference-Format style,
%% or BibLaTeX with the acmnumeric or acmauthoryear sytles, that include
%% support for advanced citation of software artefact from the
%% biblatex-software package, also separately available on CTAN.
%%
%% Look at the sample-*-biblatex.tex files for templates showcasing
%% the biblatex styles.
%%

%%
%% The majority of ACM publications use numbered citations and
%% references.  The command \citestyle{authoryear} switches to the
%% "author year" style.
%%
%% If you are preparing content for an event
%% sponsored by ACM SIGGRAPH, you must use the "author year" style of
%% citations and references.
%% Uncommenting
%% the next command will enable that style.
%%\citestyle{acmauthoryear}

%%
%% end of the preamble, start of the body of the document source.
\begin{document}

%%
%% The "title" command has an optional parameter,
%% allowing the author to define a "short title" to be used in page headers.
\title{MimiTalk: Revolutionizing Qualitative Research with Dual-Agent AI}

%%
%% The "author" command and its associated commands are used to define
%% the authors and their affiliations.
%% Of note is the shared affiliation of the first two authors, and the
%% "authornote" and "authornotemark" commands
%% used to denote shared contribution to the research.
\author{Fengming Liu}
\email{leo.liu.23@ucl.ac.uk}
\orcid{0009-0009-3881-496X}
\affiliation{%
  \institution{University College London}
  \city{London}
  \country{United Kingdom}
}

\author{Shubin Yu}
\email{yu@hec.fr}
\orcid{0000-0001-7719-3056}
\affiliation{
  \institution{HEC Paris}
  \city{Paris}
  \country{France}
}

\authornote{All authors contributed to the conceptualization and idealization of this research. F.L. led the data analysis and writing of the manuscript. S.Y. provided supervision and guidance throughout the research process. Both authors reviewed and approved the final manuscript.}


%%
%% By default, the full list of authors will be used in the page
%% headers. Often, this list is too long, and will overlap
%% other information printed in the page headers. This command allows
%% the author to define a more concise list
%% of authors' names for this purpose.
\renewcommand{\shortauthors}{Liu and Yu}

%%
%% The abstract is a short summary of the work to be presented in the
%% article.
\begin{abstract}
We present MimiTalk, a dual-agent constitutional AI framework designed for scalable and ethical conversational data collection in social science research. The framework integrates a supervisor model for strategic oversight and a conversational model for question generation. We conducted three studies: Study 1 evaluated usability with 20 participants; Study 2 compared 121 AI interviews to 1,271 human interviews from the MediaSum dataset \citep{zhuMediaSumLargescaleMedia2021} using NLP metrics and propensity score matching; Study 3 involved 10 interdisciplinary researchers conducting both human and AI interviews, followed by blind thematic analysis. Results across studies indicate that MimiTalk reduces interview anxiety, maintains conversational coherence, and outperforms human interviews in information richness, coherence, and stability. AI interviews elicit technical insights and candid views on sensitive topics, while human interviews better capture cultural and emotional nuances. These findings suggest that dual-agent constitutional AI supports effective human-AI collaboration, enabling replicable, scalable and quality-controlled qualitative research.
\end{abstract}

%%
%% The code below is generated by the tool at http://dl.acm.org/ccs.cfm.
%% Please copy and paste the code instead of the example below.
%%
\begin{CCSXML}
  <ccs2012>
    <concept>
        <concept_id>10003120.10003121.10011748</concept_id>
        <concept_desc>Human-centered computing~Empirical studies in HCI</concept_desc>
        <concept_significance>500</concept_significance>
        </concept>
    <concept>
        <concept_id>10003120.10003121.10003129</concept_id>
        <concept_desc>Human-centered computing~Interactive systems and tools</concept_desc>
        <concept_significance>500</concept_significance>
        </concept>
    <concept>
        <concept_id>10003120.10003130.10011762</concept_id>
        <concept_desc>Human-centered computing~Empirical studies in collaborative and social computing</concept_desc>
        <concept_significance>500</concept_significance>
        </concept>
    <concept>
        <concept_id>10002951.10003227.10003233</concept_id>
        <concept_desc>Information systems~Collaborative and social computing systems and tools</concept_desc>
        <concept_significance>500</concept_significance>
        </concept>
  </ccs2012>
\end{CCSXML}

\ccsdesc[500]{Human-centered computing~Empirical studies in HCI}
\ccsdesc[500]{Human-centered computing~Interactive systems and tools}
\ccsdesc[500]{Human-centered computing~Empirical studies in collaborative and social computing}
\ccsdesc[500]{Information systems~Collaborative and social computing systems and tools}



%%
%% Keywords. The author(s) should pick words that accurately describe
%% the work being presented. Separate the keywords with commas.
\keywords{Human-Computer Interaction, Human-AI Collaboration, Collaborative Research, Dual-Agent Architecture, Qualitative Research, Trust Calibration, AI-Assisted Interviews, Constitutional AI, Interview Automation, Research Methodology}
%% A "teaser" image appears between the author and affiliation
%% information and the body of the document, and typically spans the
%% page.

\begin{teaserfigure}
    \includegraphics[width=\textwidth]{figures/Teaser.jpg}
    \caption{MimiTalk: Revolutionizing Qualitative Research with Dual-Agent AI}
    \Description{MimiTalk.app}
    \label{fig:teaser}
\end{teaserfigure}

%%
%% This command processes the author and affiliation and title
%% information and builds the first part of the formatted document.
\maketitle

\section{Introduction}
The rapid advancement of generative artificial intelligence (large language models, LLMs) has fundamentally transformed how humans interact with artificial intelligence, making increasingly sophisticated conversational exchanges blur the boundaries between human and machine communication. This technological evolution presents unprecedented opportunities for automating traditionally human-intensive tasks, particularly in research contexts where conversation-based inquiry is central to data collection. However, the potential for AI systems to effectively conduct interviews—one of humanity's most fundamental social activities—remains largely unexplored in Human-Computer Interaction (HCI) research.

Traditional interview methods in qualitative research face significant challenges that directly impact user experience and research outcomes: high resource demands limit study scale, interviewer bias affects response validity, and geographical constraints restrict participant access. While video conferencing technology has addressed some distance barriers, research indicates that anonymous digital interfaces can elicit more candid responses than human interviewers, particularly regarding sensitive topics \citep{pichardRoster2020}. This finding suggests that AI-driven interview systems might offer unique advantages in certain research contexts by reducing social desirability bias and interview anxiety.

Recent empirical studies have provided evidence for AI-assisted interview methods in specific research contexts. In a pioneering study involving 395 interviews on stock market participation, researchers found that AI-assisted interviews not only maintained high data quality but achieved greater analytical depth than traditional approaches \citep{chopraConductingQualitativeInterviews2023}. Interestingly, participants in this study showed a preference for AI interviewers over human ones, suggesting that AI systems might offer advantages in certain research contexts by reducing social desirability bias and interview anxiety.

Against this background, we present MimiTalk, a novel AI-powered interview framework designed to explore how AI systems can effectively conduct interviews with human participants in qualitative research contexts. Our work addresses fundamental HCI questions about AI-powered interviewing: How can AI systems conduct interviews that maintain the depth and authenticity of traditional human-led methods while offering significant scalability advantages? What are the optimal design patterns that leverage AI capabilities in conversational generation and strategic oversight? How can trust and ethical compliance be maintained in AI-powered research frameworks?

We systematically evaluate MimiTalk through three complementary studies. Study 1 conducted a usability evaluation with 20 participants recruited through Prolific, demonstrating the platform's effectiveness in reducing interview anxiety and maintaining conversational coherence. Study 2, based on 121 English AI interviews compared against 1,271 human interviews with distinguishable roles from the MediaSum dataset \citep{zhuMediaSumLargescaleMedia2021}, demonstrates that AI-led interviews significantly outperform human interviews in information richness, semantic coherence, and result stability using NLP metrics including DeBERTa \citep{heDeBERTaDecodingenhancedBERT2020} tokenizer-based information entropy and cross-turn semantic similarity, while maintaining conversational fluency. Study 3 recruited 10 interdisciplinary bilingual researchers to conduct both human and AI interviews using consistent outlines, followed by blind thematic/content analysis of verbatim transcripts. Results show that participants generally view AI as "efficiency enhancement" rather than "creativity replacement"; trust in literature search and citation generation is low, moderate for analysis and programming, and highest for writing assistance and formatting; AI interviews can elicit technical insights comparable to expert interviews and enhance candid expression on sensitive topics like academic integrity, while human interviews excel in capturing cultural context and emotional nuances.

The main contributions of this paper are as follows:

- Conceptual contribution: Through empirical research, we further validate and confirm GenAI's role as an ``evocative object'' in interview contexts for psychological projection and meaning construction, and elucidate the necessity and mechanisms of dual-agent constitutional collaboration in high-engagement tasks; we propose and demonstrate the key tension between "scalability and emotional subtlety."
- Methodological contribution: We propose a reusable dual-agent constitutional interview paradigm and process that decouples and embeds ethical compliance, quality control, and conversational generation without sacrificing conversational naturalness.
- Empirical contribution: Based on cross-corpus comparison of 121/1,271 interviews and human-AI comparison with 10 researchers, we provide systematic evidence for AI interviews in information richness, semantic coherence, technical insights, and candid expression, while characterizing boundary conditions in cultural and emotional subtlety, offering design principles and usage recommendations for scalable qualitative research.

Next, we will review related work, detail system design and experimental methods, report quantitative and qualitative results from both studies, and finally discuss implications for HCI methodology and human-AI collaboration research.


\section{Related Work}


\subsection{AI-Driven Dialogue Systems and Human-Computer Interaction}

AI-driven dialogue systems have undergone significant evolution from simple rule-based chatbots to complex neural network architectures. Early systems relied on basic pattern matching, while the emergence of Transformer architecture \citep{vaswani2017attention} brought revolutionary advances in contextual understanding and natural conversational flow. Modern dialogue management systems can handle complex conversational contexts \citep{brabraDialogueManagementConversational2022}, providing more sophisticated solutions for interview processes and participant interactions. The development of BERT \citep{devlin2018bert} and GPT series models \citep{brown2020language} established new benchmarks for human-computer interaction, paving the way for automated interview applications. Continuous improvements in language models \citep{borgeaudImprovingLanguageModels2022} continue to enhance AI systems' conversational capabilities.

Recent empirical studies have provided evidence for AI-assisted interviews in large-scale research. In a study involving 395 interviews on stock market participation, researchers found that AI-assisted interviews not only maintained high data quality but achieved deeper analytical results than traditional methods \citep{chopraConductingQualitativeInterviews2023}. Interestingly, participants in this study showed a preference for AI interviewers over human ones, challenging traditional assumptions about qualitative research methods.

The success of AI in interviews has extended beyond financial research. In recruitment contexts, AI systems demonstrate sophisticated conversational capabilities that enhance candidate engagement \citep{ahmadFutureRecruitmentUsing2024}. In healthcare, AI improves clinical decision-making through optimized patient interactions \citep{alowaisRevolutionizingHealthcareRole2023, krishnanArtificialIntelligenceClinical2023}. The financial services industry utilizes AI for predictive analysis and customer interactions \citep{bahooArtificialIntelligenceFinance2024}, while education leverages AI to create more personalized learning experiences \citep{sapciArtificialIntelligenceEducation2020}. These cross-domain applications demonstrate that AI systems can support diverse forms of human-computer interaction in specific contexts.

Field studies have also revealed interesting dynamics in AI interview applications. For example, research on asynchronous video interviews (AI-AVI) shows that when systems integrate perceptibility and transparency features, participants' cognitive trust increases \citep{suenHung2023}. The development of AI conversational interview systems \citep{wuttkeAIConversationalInterviewing2024} further expands the possibilities of automated interviews. In a large-scale study involving 710 simulated interviews, machine learning models successfully predicted personality traits (R² = 0.32) and interview performance (R² = 0.44), with accuracy exceeding self-report indicators \citep{koutsoumpis2024}. However, such evaluation systems still face challenges in non-evaluative applications like social science research.

The integration of AI systems in research contexts raises important questions about human-AI collaboration and trust calibration. Research shows that even in high-scoring tasks, users remain cautious about completely delegating citations, fact-checking, or interpretive judgments to AI, with a general "trust ceiling" of approximately 80\% \citep{openaiGPT4TechnicalReport2024}. This skepticism reflects ongoing concerns in the academic community about hallucination phenomena and limitations in AI contextual understanding. Research on retrieval-augmented generation \citep{shusterRetrievalAugmentationReduces2021} indicates that anchoring AI responses to factual information can reduce hallucination problems. Meanwhile, efforts to improve fairness in machine learning systems \citep{holsteinImprovingFairnessMachine2019} address important ethical considerations in AI applications.

In the field of AI-assisted qualitative data analysis, recent research has further validated the potential of AI tools in social science research. Akman et al. \citep{akmanHumanResearcherVs2025} compared human researchers with AI-supported qualitative data analysis and found that hybrid prompt design can significantly enhance the efficiency and accuracy of qualitative data analysis, providing deeper insights and more structured outputs. The study employed open, axial, and selective coding methods, demonstrating AI's effectiveness in programming education research while emphasizing the importance of carefully designed prompts and human oversight in reducing potential biases and errors in AI analysis. Hamilton et al. \citep{hamiltonExploringUseAI2023} compared human and AI-generated qualitative analysis in guaranteed income data research, finding both similarities and differences in theme identification, with human coders identifying some themes that ChatGPT did not, and vice versa. This finding suggests that AI can serve as a complementary tool for complex human-centered tasks, providing additional support for research tasks. Li et al. \citep{liComparingGPT4Human2024} compared GPT-4 with human researchers in healthcare data analysis, finding that GPT-4 can effectively identify key themes, showing moderate agreement with human analysis ($\kappa=0.401$), and while human analysis provided richer subtheme diversity, AI's consistency suggests its potential as a complementary tool for qualitative research.

\subsection{Dual-AI Agent as "evocative object"}
The question of whether artificial intelligence should be regarded as a tool for social science research or as a "conversationalist" with its own dialogical status remains highly controversial. Contemporary scholars often reduce the dialogue and cognitive experiences created by AI systems to purely instrumental use, essentially viewing them as deterministic, algorithmic mechanical processes. This perspective can be analogized to Searle's "Chinese Room" thought experiment \citep{searleMindsBrainsPrograms1980}. In this scenario, a person who doesn't understand Chinese is locked in a room, mechanically matching and outputting Chinese symbols based solely on an operation manual. To external observers, it appears that this person understands Chinese, but in fact, they have no real understanding. If we believe that AI systems lack true understanding or consciousness, merely manipulating symbols according to preset rules without grasping their meaning, then we are operating within the framework of the "Chinese Room."

However, despite the "Chinese Room" emphasizing AI's lack of true understanding, Dennett's "intentional stance" theory \citep{dennettIntentionalStance1998} provides a perspective more aligned with our research context. Dennett argues that when we adopt an "intentional stance" toward a system, we interpret its behavior as that of a rational agent, assuming it has beliefs, desires, and intentions. This strategy is often very effective in predicting and understanding system behavior. In our research context, this theoretical perspective helps explain why participants can engage in meaningful dialogue with AI even when they know that AI lacks true consciousness.

Furthermore, subjectivity need not be understood as an entity inherent in consciousness, but can be viewed as a position within the symbolic order. Lacan's psychoanalytic theory \citep{lacanEcritsFirstComplete2006} indicates that subjectivity is constructed through language and symbolic systems, where the subject does not exist as a transcendental autonomous entity but occupies specific symbolic positions. Žižek further argues that the core of ideological operation lies in the subject's identification with and projection onto symbolic authority, maintaining belief even when the subject knows of its fictionality \citep{zizekSublimeObjectIdeology2002}. These theories reveal a persistent impression among AI scientists and algorithm engineers: that there is a fundamental difference between human intelligence and the mechanical, programmatic nature of machines. However, in human-computer interaction contexts, if humans cannot infer understanding from AI's input and output, then they also cannot determine whether other humans truly possess such understanding. The only difference is that humans can base empathy on analogies to their own intelligence, and this analogy cannot naturally extend to AI.

Moreover, the advanced natural language processing capabilities of contemporary AI systems make their outputs no longer simply predictable. This is fundamentally different from text correction algorithms like Grammarly or Word: when humans interact with AI, they often expect results that exceed their own predictions, even hoping that AI can correct or expand their own thinking.

This phenomenon forms an interesting contrast with the findings of Luger et al. \citep{lugerHavingReallyBad2016}. Their research shows that users often view conversational agents as anthropomorphic assistants with subjectivity, but actual experiences are often disappointing, revealing gaps between expectations and reality.

This also explains why participants, even knowing its artificial nature, tend to attribute real subjectivity to our collaborative framework. This phenomenon reflects the psychological mechanism of humans' tendency to project social rules and expectations onto computers, as revealed by media equation research \citep{nassMachinesMindlessnessSocial2000}. We believe that MimiTalk's collaborative framework functionally resembles what Turkle describes as an "evocative object" \citep{turkleAloneTogetherWhy2011}—a technological entity capable of stimulating psychological projection and meaning construction. Due to its lack of human judgment, emotional responses, and social expectations, it instead creates a psychological space where participants can express themselves more authentically, while maintaining research quality under the premise of human researchers providing strategic supervision and contextual understanding.

\subsection{Applications of Human-Computer Collaboration in Research Contexts}

The evolution of interview methodology has been significantly influenced by technological advances, particularly in the field of human-computer interaction. Traditional interview methods are generally divided into three categories: structured interviews with fixed questions, semi-structured interviews allowing some flexibility, and unstructured interviews emphasizing open dialogue \citep{cuevas2024}. These methods have gradually evolved from simple conversations to complex approaches integrating digital interfaces and multimodal interactions.

Recent HCI research has explored the impact of digital interfaces on interview dynamics. Studies show that compared to human-conducted interviews, anonymous digital interfaces can elicit more candid responses from participants when dealing with sensitive topics \citep{pichardRoster2020}. This finding suggests that AI-driven interview systems may have unique advantages in certain research contexts, reducing social desirability bias and interview anxiety.

The theoretical framework of digital interviews has evolved from standardized protocols in the early 20th century to contemporary approaches with greater cultural sensitivity \citep{brinkmann2018}. The formulation of questions plays a crucial role in knowledge acquisition, similar to how different lenses affect photographic effects \citep{briggs2003questionformulation}. Various interview methods serve different purposes: structured interviews ensure standardization, semi-structured interviews balance flexibility with systematic inquiry, while unstructured interviews encourage free narration \citep{kvale2021, seidman2019}. The concept of "active interviewing" \citep{holsteinActiveInterview1995} further deepens our understanding of how interviewers actively participate to gain deeper insights. Research on computer-automated interviews \citep{pickardUsingComputerAutomated2020} has also demonstrated the feasibility and effectiveness of automated systems in conducting structured interviews. Comparative studies of online and offline interviews \citep{shapkaOnlineInpersonInterviews2016} reveal trade-offs between different interview modes.

However, traditional interview methods still face some challenges that AI-driven systems may provide solutions for: for example, high resource consumption severely limits research scale \citep{seidman2019}, interviewer bias affects response validity \citep{kvale2021}, and geographical constraints hinder participant accessibility. Although video conferencing technology has alleviated distance barriers, research also indicates that it may weaken the quality of interpersonal interactions \citep{brinkmann2018}.

\section{AI-Based Interview System Design}

\subsection{Platform Architecture}
We built the application platform Mimitalk.app for the MimiTalk framework, providing it for researchers and interviewees to use. MimiTalk.app adopts a modular architecture consisting of three core components: frontend interview module, backend API service, and AI inference module. The frontend is implemented based on React.js and Material UI components, providing intuitive interaction interfaces and smooth user experience for interviewers and interviewees. The backend API service is implemented based on FastAPI, providing RESTful interfaces for interview management and real-time communication based on WebSocket.

We adopted an approach where the AI interviewer actively initiates conversations, followed by users actively starting recording and using the Whisper speech model for text conversion during interviews. We believe that using AI for interviews is a novel and unique qualitative data collection method, rather than a pure simulation and replacement of human interviews. To control the immersion generated by task avatars, the MimiTalk.app interview interface does not include human-like avatars and allows users to view interviewer questions and real-time interview transcripts in text form. To reduce users' cognitive burden, the MimiTalk.app interview interface adopts a simple minimalist design, retaining only the recording button for starting responses and the button for ending interviews.

\begin{figure}[t]
  \centering
  \begin{subfigure}[b]{0.48\textwidth}
    \centering
    \includegraphics[width=\linewidth]{figures/Screenshot.png}
    \caption{Screenshot of the MimiTalk.app interview interface. The interface adopts a minimalist design, retaining only core function buttons such as recording and ending the interview, avoiding anthropomorphic avatars and highlighting the conversation content.}
    \label{fig:screenshot1}
  \end{subfigure}
  \hfill
  \begin{subfigure}[b]{0.48\textwidth}
    \centering
    \includegraphics[width=\linewidth]{figures/ai_interview_architecture_chi.png}
    \caption{MimiTalk dual-agent collaborative interview system architecture. The process covers interview initialization, question generation, and response handling. The supervisor model is responsible for strategic control, while the main model generates natural conversation, ensuring compliance and fluency.}
    \label{fig:ai_logic_flow}
  \end{subfigure}
  \caption{Side-by-side display of the MimiTalk.app interview interface and system architecture: (a) The interview interface features a minimalist design to avoid confounding variables and focus on interview content; (b) The dual-agent collaborative architecture integrates constitutional principles and real-time context analysis.}
\end{figure}

\subsection{Dual-Agent Constitutional AI Architecture}
This system implements a dual-agent architecture for constitutional collaborative GenAI, with supervisor and conversational models working together to achieve controllable and scalable research interviews while ensuring ethical compliance and conversational naturalness.

\textbf{System Architecture Core}  
The core of the system lies in embedding constitutional AI principles into the dual-agent collaborative architecture, achieving organic integration of strategic supervision and natural conversation through specialized division of labor. The system adopts a modular design that decouples supervision logic from conversational generation.

\textbf{Supervisor Agent}  
The supervisor agent uses the Claude Sonnet 4.0 model, responsible for strategic-level interview quality control and ethical compliance. It continuously monitors conversational development, assesses compliance with predefined objectives, and provides real-time directional suggestions to maintain interview depth and effectiveness while ensuring ethical compliance.

\textbf{Core Multi-Model Asynchronous Architecture}  
The system implements an asynchronous architecture where the responder agent uses GPT-5 or Claude Haiku models for natural conversational generation, with intelligent model selection and background supervision analysis:

The complete implementation of the core multi-model asynchronous architecture, including the intelligent model selection and background supervision analysis functions, is provided in Appendix A.

This asynchronous architecture enables natural conversational generation with dynamic response adaptation, maintaining conversational continuity and relevance while ensuring all generated content complies with constitutional AI principles.

Figure \ref{fig:ai_logic_flow} shows the dual-model architecture for AI interview question generation. The process covers input processing, supervisor model analysis and strategic guidance generation, prompt construction and conversational generation steps.

\section{Study 1: Usability Evaluation of Mimitalk.app}

To validate the effectiveness and user experience of our AI-powered interview platform, Study 1 conducted a comprehensive usability evaluation involving 20 participants recruited through Prolific, an online research participant recruitment platform. This study serves as a foundational assessment of the platform's technical capabilities, user interface design, and overall interview experience before proceeding to more complex comparative analyses.

The participants represented diverse demographic backgrounds and geographical locations, ensuring a representative sample for our usability study. The participant pool included individuals aged 18 to 70 years (median age early 30s), with a relatively balanced gender distribution (14 female, 12 male among known participants). Educational backgrounds varied from high school/A-Levels to Master's degrees, with Bachelor's degrees being the most common qualification. The occupational profile was notably diverse, including students, professionals (such as HR advisors, financial analysts, and teachers), healthcare workers, administrative staff, homemakers, and retired individuals. Participants were predominantly from the United Kingdom, and most reported prior interview experience, primarily in job-seeking contexts, across different formats including face-to-face, online, and telephone interviews. More information of the participants in Study 1 can be found in Appendix \ref{sec:usability_demographics}.

The interview outline was structured to assess both the platform's technical capabilities and user experience. The outline can be found in Appendix \ref{sec:usability_test_outline}.

The analysis revealed that AI-generated questions showed significant advantages over static approaches through their ability to maintain conversational coherence while exploring participant expertise. However, limitations were observed including occasional redundancy, context drift in longer conversations, and cultural sensitivity gaps where AI did not always recognize nuances in academic hierarchies. The platform's ability to generate follow-up questions based on previous responses was particularly effective, with 14 of the participants noting that the questions felt contextually appropriate and engaging. 

\begin{table}[h]
\caption{Key Themes from Usability Study Qualitative Analysis}
\label{tab:usability_themes}
\small
\begin{tabular}{@{}p{4cm}|p{8cm}@{}}
\toprule
\textbf{Theme} & \textbf{Description} \\
\midrule
User Experience & Platform usability, interface interaction, response mechanisms \\
Psychological Impact & Anxiety levels, comfort, pressure management \\
Technical Performance & Response time, accuracy, system reliability \\
Comparative Elements & Differences from traditional interviews \\
Limitations & Areas for improvement and development \\
\bottomrule
\end{tabular}
\end{table}

The usability evaluation revealed several key insights about the AI-powered interview platform's effectiveness and user experience. The system indicated 41 people clicked the interview link, with 20 participants completing the full interview process, representing a $48.8\%$ completion rate that demonstrates reasonable user engagement.

Participant feedback highlighted several strengths, including natural conversation flow, comfortable interview environment, and clear and logical question progression. The interface received particular praise for its intuitive design and straightforward navigation. Many participants emphasized the platform's quick response times and accurate speech recognition capabilities, which contributed to a smooth interview process.

The psychological aspects of the AI-driven interview process emerged as a significant theme. Participants consistently reported lower anxiety levels compared to traditional interviews, citing the absence of immediate human judgment as a key factor. The ability to pace their responses and focus on content rather than physical presentation was particularly valued by participants who reported previous interview anxiety.

When comparing the platform to traditional interview experiences, participants also noted several distinctive advantages. The consistency in question delivery and standardized evaluation approach was frequently mentioned as a positive feature. Participants particularly appreciated the elimination of appearance-based bias and reduced pressure regarding body language, allowing them to focus more on their responses' substance.

However, limitations were identified including the lack of real-time response transcription, limited ability to build personal rapport compared to human-led interviews, and occasional challenges with natural language understanding. The overall sentiment was balanced, with participants endorsing the platform's effectiveness as a preliminary screening tool and suggesting its value as a complementary tool.
\section{Methodology}


\subsection{Study 2: Quantitative Evaluation of AI-powered Interviews}

\textit{Data Collection and Privacy Considerations.} All AI interview data for this study were collected through the MimiTalk.app platform, with research permissions included in the platform's terms of service. Due to user privacy protection and ethical considerations, we cannot share the raw interview data with third parties. However, future research may consider collecting additional AI interview data separately for replication studies or larger-scale investigations and mechanism studies.

To deeply evaluate the quality of AI-driven interviews, we designed a multi-dimensional analysis framework, combining natural language processing metrics such as information entropy and cross-turn semantic similarity, based on DeBERTa tokenizer \citep{heDeBERTaDecodingenhancedBERT2020} and DeBERTa v3 \citep{heDeBERTaV3ImprovingDeBERTa2021} for analysis. To ensure the robustness of results, we used early BERT models \citep{devlinBERTPretrainingDeep2018} as a control benchmark. Under the MimiTalk Framework, we conducted bilingual corpus comparative analysis of 121 English semi-structured AI interviews (each with transcribed text exceeding 3000 characters) against human-guided interview texts from the MediaSum dataset \citep{zhuMediaSumLargescaleMedia2021}, and tested for significant differences through t-tests. 

\textit{1. Information Entropy}

Information entropy, as a core metric for measuring text complexity and information content, has wide application value in natural language processing and social science fields. Higher entropy reflects greater informational complexity and potentially enhanced interview efficiency \citep{dahlQuantifyingInformationContent2010, friesnerInformationEntropyScale2021}. This approach has proven particularly effective in research contexts where measuring information diversity and content richness is essential for evaluating response quality \citep{homolaMeasureSurveyMode2016, keykhaAdvantagesChallengesElectronic2025}. For example, Dahl and Østerås \citep{dahlQuantifyingInformationContent2010} applied Shannon entropy to survey data, proposing the concept of information efficiency, which is the ratio of empirical entropy to maximum achievable entropy.

We use the DeBERTa base model \citep{heDeBERTaDecodingenhancedBERT2020} for subword tokenization, which provides superior performance in handling mixed-language content compared to traditional word-based approaches.

We calculate the word-level information entropy using the Shannon entropy formula:
\begin{equation}
  H(X) = -\sum_{i=1}^{n} p(x_i) \log_2 p(x_i)
\end{equation}
where $H(X)$ represents the information entropy, $p(x_i)$ is the probability of the $i$-th token, and $\log_2$ denotes the logarithm to base 2.

For each token $i$:
\begin{equation}
  p(x_i) = \frac{\text{count}(x_i)}{\text{total\_tokens}}
\end{equation}

To eliminate the influence of text length on entropy values, we normalize all transcripts to the minimum length across all interviews using truncation. Specifically, if $L_{\min}$ is the minimum token count across all interviews, we truncate longer texts to $L_{\min}$ tokens while preserving the original order of tokens.

The higher resulting entropy values indicate greater lexical diversity and information richness. An entropy of 0 would indicate a text containing only one repeated token, while higher entropy values suggest more diverse vocabulary usage.

To measure the depth and diversity of responses elicited by AI-conducted interviews, we calculate the entropy values separately for interviewer text, interviewee text, and the overall text. Using the same tokenizer, we also calculate sentence length as an auxiliary reference metric. We calculate descriptive statistics (mean, standard deviation, minimum, maximum) for entropy and sentence length across AI-conducted and human-conducted interviews, and assess statistical significance using independent samples t-tests with Welch's correction for unequal variances.


\textit{2. Semantic Similarity}

Semantic similarity measures the degree of conceptual relatedness between textual units, quantifying how closely different utterances align in meaning rather than just lexical overlap. This metric is crucial for evaluating interview quality because it captures conversational coherence—the extent to which participants maintain thematic consistency and engage meaningfully with each other's contributions. High semantic similarity indicates coherent dialogue flow, where questions build upon previous responses and participants stay on-topic, while low similarity may suggest fragmented or disjointed conversations.

In interview contexts, semantic similarity serves three key purposes: measuring how consistently individual speakers maintain thematic focus within their own contributions, indicating depth of engagement; evaluating how well interviewer questions relate to interviewee responses, reflecting conversational quality and mutual understanding; providing a holistic assessment of interview structure and flow. These measures are particularly valuable for comparing AI and human interviewers because they reveal whether AI systems can maintain the natural conversational dynamics that characterize effective human-led interviews.

In social science and human-computer interaction research, the potential of BERT and its derivative models has been widely validated. Fang et al. \citep{fangEvaluatingConstructValidity2022} constructed a validity framework, proving the superiority of BERT embedding technology in representing survey questions, with stronger semantic capture capabilities compared to traditional methods, providing technical assurance for semantic analysis in AI interviews. Grandeit et al. \citep{grandeitUsingBERTQualitative2020} applied BERT to qualitative analysis of psychosocial online counseling, developing an automated coding system with over 50 categories, significantly reducing manual annotation workload. This indicates that AI interviews have high efficiency and scalability in processing large-scale text data. Wankmüller's \citep{wankmullerIntroductionNeuralTransfer2021} research further compared Transformer models with traditional machine learning algorithms, finding that BERT can maintain high prediction accuracy even with limited data, which is particularly crucial for AI interview scenarios. Mostafavi et al. \citep{mostafaviContextualEmbeddingsSociological2025} developed the BERTNN method, extending emotion dictionaries through BERT's contextual embeddings, providing efficient tools for emotion analysis in sociological research and greatly reducing time costs of traditional surveys. Additionally, Aral and Dhillon's \citep{aralWhatExactlyNovelty2023} research explored the sources of information novelty in network structures, pointing out the advantages of weak connections in transmitting unique information, providing theoretical inspiration for understanding AI interviews' role in information transmission. Keykha et al. \citep{keykhaAdvantagesChallengesElectronic2025} analyzed the advantages and disadvantages of electronic examinations through Shannon entropy, revealing the potential and challenges of technology-driven tools in educational assessment, providing reference for automated assessment in AI interviews.

We employ pre-trained transformer models for sentence-level semantic similarity analysis. Specifically, we use the DeBERTa v3 model \citep{heDeBERTaV3ImprovingDeBERTa2021} and BERT base model \citep{devlinBERTPretrainingDeep2018} to generate sentence embeddings, then calculate cosine similarity between sentence pairs.

For sentence segmentation, we split text using punctuation markers ($[.!?]+$) and speaker indicators. Each sentence $s_i$ is tokenized and embedded using the transformer model:
\begin{equation}
\mathbf{e}_i = \text{Model}(s_i) \in \mathbb{R}^{768}
\end{equation}

where $\mathbf{e}_i$ is the 768-dimensional sentence embedding obtained by mean pooling over the last hidden states of the transformer model.

In the semantic similarity analysis, we separately calculated the cosine similarity of sentences within the interviewer, within the interviewee, and between the two speakers. Specifically, let $\mathbf{e}_i^{\text{int}}$ denote the embedding vector of the $i$-th interviewer sentence, $N_{\text{int}}$ the number of interviewer sentences, $\mathbf{e}_j^{\text{intv}}$ the embedding vector of the $j$-th interviewee sentence, and $N_{\text{intv}}$ the number of interviewee sentences. The formulas for the three types of similarity are as follows:
\begin{equation}
\text{Sim}_{\text{int}} = \frac{1}{N_{\text{int}}(N_{\text{int}}-1)/2} \sum_{i<j} \frac{\mathbf{e}_i^{\text{int}} \cdot \mathbf{e}_j^{\text{int}}}{\|\mathbf{e}_i^{\text{int}}\| \|\mathbf{e}_j^{\text{int}}\|}
\end{equation}
\begin{equation}
\text{Sim}_{\text{intv}} = \frac{1}{N_{\text{intv}}(N_{\text{intv}}-1)/2} \sum_{i<j} \frac{\mathbf{e}_i^{\text{intv}} \cdot \mathbf{e}_j^{\text{intv}}}{\|\mathbf{e}_i^{\text{intv}}\| \|\mathbf{e}_j^{\text{intv}}\|}
\end{equation}
\begin{equation}
\text{Sim}_{\text{cross}} = \frac{1}{N_{\text{int}} N_{\text{intv}}} \sum_{i,j} \frac{\mathbf{e}_i^{\text{int}} \cdot \mathbf{e}_j^{\text{intv}}}{\|\mathbf{e}_i^{\text{int}}\| \|\mathbf{e}_j^{\text{intv}}\|}
\end{equation}

Here, $\cdot$ denotes the dot product, and $\|\cdot\|$ denotes the L2 norm.

The cosine similarity ranges from -1 to 1, where 1 indicates identical semantic content, 0 represents orthogonal (unrelated) content, and -1 denotes opposite semantic content.

For robustness testing, we compare results across multiple transformer models (BERT and DeBERTa v3) to ensure consistency of findings.

We calculate descriptive statistics (mean, standard deviation, minimum, maximum, median) for each similarity type across AI-conducted and human-conducted interviews, and assess statistical significance using independent samples t-tests with Welch's correction for unequal variances.

\textit{3. Causal identification}

We employ Propensity Score Matching (PSM) to identify the causal effects of AI-conducted interviews on linguistic characteristics. This quasi-experimental approach addresses potential selection bias by balancing observable confounders between AI and human interview groups.

We construct covariates using interview-level characteristics that affect interview type selection but are not directly influenced by AI technology:
\begin{equation}
X_i = [\text{total\_tokens}_i, \text{total\_sentences}_i, \text{PC1}_i, \text{PC2}_i, \text{PC3}_i]
\end{equation}

where $\text{total\_tokens}_i$ and $\text{total\_sentences}_i$ represent interview length and dialogue turns, and $\text{PC1}_i, \text{PC2}_i, \text{PC3}_i$ are principal components from topic embedding analysis. We apply log transformation, polynomial features, and quantile transformation to improve covariate distributions.

We estimate propensity scores using logistic regression:
\begin{equation}
P(T_i = 1 | X_i) = \frac{\exp(\beta_0 + \beta_1 X_{i1} + ... + \beta_k X_{ik})}{1 + \exp(\beta_0 + \beta_1 X_{i1} + ... + \beta_k X_{ik})}
\end{equation}

where $T_i = 1$ indicates AI-conducted interview and $T_i = 0$ indicates human-conducted interview.

We employ Kernel Matching with Gaussian kernel function and adaptive caliper $c = 0.2 \times \text{SD}(PS)$:
\begin{equation}
K(d_{ij}) = \exp\left(-\frac{d_{ij}^2}{2h^2}\right)
\end{equation}

where $d_{ij} = |PS_i - PS_j|$ is the propensity score distance and $h = 0.1$ is the bandwidth parameter.

On the matched sample, we estimate Average Treatment Effect (ATE) using Ordinary Least Squares:
\begin{equation}
Y_i = \alpha + \beta T_i + \gamma_1 X_{i1} + ... + \gamma_k X_{ik} + \epsilon_i
\end{equation}

where $Y_i$ is the outcome variable (entropy, sentence length, or semantic similarity), $T_i$ is the treatment indicator, and $\beta$ is the ATE estimate.

We conduct placebo tests using digit ratio and sentence count variables theoretically unaffected by AI interview technology, and perform alternative metric testing using DeBERTa v3 embeddings for semantic similarity. Our causal identification relies on unconfoundedness ($Y(1), Y(0) \perp T | X$), overlap ($0 < P(T = 1 | X) < 1$), and SUTVA (no interference between units). We assess covariate balance using Standardized Mean Difference (SMD) and report ATE estimates with standard errors and 95\% confidence intervals.

We address potential endogeneity issues through comprehensive length normalization mechanisms. For semantic similarity analysis, we implement sentence-level normalization, embedding dimension normalization, and similarity calculation normalization, ensuring that \texttt{total\_tokens} and \texttt{total\_sentences} as covariates do not introduce endogeneity problems. For information entropy analysis, we apply text length normalization, text truncation standardization, and information entropy calculation normalization, where entropy values are computed based on standardized text lengths and all interview texts are truncated to the same minimum length. These normalization procedures eliminate the potential endogeneity concerns that could arise from text length variations affecting our outcome variables.

\subsection{Study 3: Comparative Analysis of Human-led vs. AI-conducted Interviews}

To further validate the deep insight capabilities of AI systems, we evaluated whether AI systems can achieve or even surpass the depth and authenticity of traditional human interviewers by comparing human-led and AI-led interview approaches. We chose a more precise controlled variable method, with this design philosophy based on a comparative case study methodological framework that has been proven effective in evaluating different data analysis methods \citep{akmanHumanResearcherVs2025}. We therefore conducted human interviews and AI semi-structured interviews under the MimiTalk framework using consistent outlines, followed by blind thematic/content analysis of verbatim transcripts.

We recruited 10 academic researchers across diverse disciplinary backgrounds, ensuring methodological rigor through representational diversity. Participants included PhD students, lecturers, postdoctoral researchers, and professors with 0.5 to 10 years of research experience, all bilingual in Chinese and English. Disciplinary backgrounds spanned philosophy (including Marxist theory), microbiology, computer science (vehicular networking and model optimization), economics (economic history), strategic management (innovation and human-AI collaboration), atmospheric science, speech recognition, and image processing. Detailed demographic and research background information is provided in Appendix B.

Semi-structured interviews, as an important method in qualitative research, have wide application in social science research. This method combines the standardization advantages of structured interviews with the flexibility of open interviews, allowing researchers to conduct in-depth follow-up questions based on respondents' answers while maintaining focus on core themes \citep{brinkmannQualitativeInterviewing2013}. In AI interview systems, semi-structured design is particularly important because it provides sufficient guidance framework for AI systems while retaining the ability to adjust questions based on conversational dynamics. We conducted blind thematic/content analysis of verbatim transcripts, with reviewers (unaware of interview sources) using reflective thematic analysis methods to evaluate the thematic depth, emotional subtlety, and information richness of texts. This blind review method has been proven to effectively reduce subjective bias and is widely used in medical and social science qualitative research. For example, Sakaguchi et al.'s \citep{sakaguchiEvaluatingChatGPTQualitative2025} study comparing ChatGPT with humans in Japanese clinical contexts showed that AI achieved over 80\% agreement on descriptive themes but only about 30\% on cultural emotional themes, inspiring us to pay special attention to AI's boundaries in handling sensitive topics. Similarly, Prescott et al. \citep{prescottComparingEfficacyEfficiency2024} compared humans and generative AI in qualitative thematic analysis, finding that AI achieved 71\% thematic consistency in inductive analysis and significantly reduced analysis time (20 minutes vs. 567 minutes), providing empirical foundation for emphasizing the necessity of hybrid modes. Gibson and Beattie \citep{gibsonMoreLessHuman2024} further pointed out that AI struggles to replicate the richness of human experience in online qualitative research, emphasizing the role of theoretical frameworks in identifying AI-generated data, which aligns with our blind review method.

In interviews conducted under the Mimitalk Framework, we introduced more accurate outlines and protocols, informing the AI system in advance about the content researchers want to explore and potential issues they might encounter through examples and guiding words, combined with the AI system's self-adjustment capabilities, making interviews more natural and in-depth. This prompt engineering method has been proven to significantly improve analysis efficiency and accuracy in AI-assisted qualitative analysis \citep{akmanHumanResearcherVs2025}. Our prompts designed four themes, including: AI ethics and responsibility, AI potential risks and challenges, AI practical applications and effects, and AI future development and prospects.

Interdisciplinary research design has significant importance in evaluating the effectiveness of AI systems in different field applications. By recruiting researchers from different disciplinary backgrounds, we were able to assess the adaptability and effectiveness of AI interview systems in diverse research contexts. This design method not only enhanced the external validity of research results but also provided a comprehensive perspective for understanding AI system performance in different professional fields \citep{stokolsScienceTransdisciplinaryAction2006}. The participation of bilingual researchers further ensured the applicability of research results in cross-cultural contexts, which is crucial for evaluating the potential of AI systems in globalized research environments.

English Prompts (translated): {Brief understanding: research field, years of work experience, working language
1. If you were to rate the academic research value of AI in your field on a scale of 1-5 (1=completely useless, 5=transformative), what would your score be? Could you briefly explain the reasons? Please explain from perspectives such as research efficiency, innovation, and quality of outcomes.
2. Compared to traditional methods, has AI assistance changed your time allocation patterns? In which specific work stages or steps have you felt changes in time allocation: for example, in literature search, data analysis, writing drafts, etc., have there been obvious time reallocations?
3. From 0\% (completely distrust) to 100\% (completely trust), what is your reliability rating for AI-generated academic content? Please explain this value with specific scenarios [If below 80\%, ask: How would you verify AI-generated content when reliability is insufficient?]
4. In literature review/experimental design/data collection/paper writing/outcome review, which stage do you think AI is most suitable to intervene in? Which is least suitable? Please provide examples. Are there any unlisted stages where you think AI might be applicable?
5. What are your expectations and concerns about the future of AI in academic research? What opportunities might it bring? What risks should we be vigilant about? (Please provide examples)
6. In your research process, have AI tools ever had a significant impact on your conclusions or research direction? Please provide examples. Have you ever adjusted research methods, modified hypotheses, or discovered new research questions because of AI suggestions?}

\section{Results}

\subsection{Study 2}
\subsubsection{Descriptive Statistics}
Our comprehensive analysis of 121 AI interviews compared to 1,271 human interviews from the MediaSum dataset reveals significant differences across three key dimensions: information entropy, token usage patterns, and semantic similarity.

\textbf{Information Entropy Analysis}

Information entropy analysis demonstrates that AI-conducted interviews consistently exhibit higher linguistic diversity across all measured categories. Overall transcript entropy shows AI interviews achieving significantly higher values (7.703 ± 0.399) compared to human interviews (7.273 ± 0.395), representing a 5.9\% increase in vocabulary diversity (Figure \ref{fig:entropy_overall}). This pattern extends to both speaker roles: AI interviewee responses demonstrate higher entropy (7.098 ± 0.791) than human counterparts (6.907 ± 0.446), while AI interviewer text shows the most pronounced difference (7.325 ± 0.512 vs. 6.762 ± 0.401), representing an 8.3\% increase in question diversity (Figures \ref{fig:entropy_interviewee} and \ref{fig:entropy_interviewer}). The comprehensive comparison across all categories confirms AI interviews' superior linguistic richness (Figure \ref{fig:comprehensive_entropy}).

\begin{figure}[h]
  \centering
  \begin{subfigure}[b]{0.32\textwidth}
    \centering
    \includegraphics[width=\textwidth]{figures/overall_transcript_entropy.png}
    \caption{Overall transcript entropy}
    \label{fig:entropy_overall}
  \end{subfigure}
  \hfill
  \begin{subfigure}[b]{0.32\textwidth}
    \centering
    \includegraphics[width=\textwidth]{figures/interviewee_response_entropy.png}
    \caption{Interviewee response entropy}
    \label{fig:entropy_interviewee}
  \end{subfigure}
  \hfill
  \begin{subfigure}[b]{0.32\textwidth}
    \centering
    \includegraphics[width=\textwidth]{figures/interviewer_text_entropy.png}
    \caption{Interviewer text entropy}
    \label{fig:entropy_interviewer}
  \end{subfigure}
  \caption{Information entropy distributions comparing AI and human interviews. (a) Overall transcript entropy: AI interviews demonstrate significantly higher linguistic diversity (7.703 ± 0.399) compared to human interviews (7.273 ± 0.395), representing a 5.9\% increase in vocabulary richness. (b) Interviewee response entropy: AI interviewees exhibit higher entropy (7.098 ± 0.791) than human interviewees (6.907 ± 0.446), indicating more diverse vocabulary usage in responses. (c) Interviewer text entropy: AI interviewers show substantially higher entropy (7.325 ± 0.512) compared to human interviewers (6.762 ± 0.401), representing an 8.3\% increase in question diversity.}
\end{figure}

\begin{figure}[h]
  \centering
  \begin{subfigure}[b]{0.32\textwidth}
    \centering
    \includegraphics[width=\textwidth]{figures/comprehensive_entropy_comparison.png}
    \caption{Comprehensive entropy comparison}
    \label{fig:comprehensive_entropy}
  \end{subfigure}
  \hfill
  \begin{subfigure}[b]{0.32\textwidth}
    \centering
    \includegraphics[width=\textwidth]{figures/token_count_kde_distribution.png}
    \caption{Token count KDE}
    \label{fig:token_kde}
  \end{subfigure}
  \hfill
  \begin{subfigure}[b]{0.32\textwidth}
    \centering
    \includegraphics[width=\textwidth]{figures/similarity_comparison.png}
    \caption{Similarity comparison}
    \label{fig:similarity_comparison}
  \end{subfigure}
  \caption{Comprehensive analysis summaries. (a) Comprehensive information entropy comparison across all categories: The violin plots demonstrate that AI interviews consistently achieve higher linguistic diversity across overall transcripts, interviewer questions, and interviewee responses, with tighter distributions indicating more reliable performance. (b) Kernel density estimation of token count distributions: The density plots reveal distinct patterns where AI interviews show broader, more right-skewed distributions for both interviewers and interviewees, indicating greater variability in response lengths compared to human interviews. (c) Comprehensive semantic similarity comparison across all categories: The violin plots demonstrate that AI interviews consistently outperform human interviews in semantic coherence across interviewer internal, interviewee internal, and cross-speaker similarity measures.}
\end{figure}

\textbf{Token Usage Patterns}

Token count analysis reveals distinct patterns in response length and variability between AI and human interviews. AI interviewees produce significantly longer responses (24.7 ± 22.2 tokens per sentence) compared to human interviewees (18.8 ± 15.5 tokens), representing a 31.4\% increase in response length (Figure \ref{fig:token_interviewee}). AI interviewers also generate slightly longer questions (16.5 ± 11.6 tokens) than human interviewers (14.5 ± 11.3 tokens), showing a 13.8\% increase (Figure \ref{fig:token_interviewer}). Notably, AI interviews exhibit greater variability in response lengths, particularly among interviewees (standard deviation: 22.2 vs. 15.5 tokens), suggesting more diverse response strategies (Figures \ref{fig:token_overall} and \ref{fig:token_kde}).

\begin{figure}[h]
  \centering
  \begin{subfigure}[b]{0.32\textwidth}
    \centering
    \includegraphics[width=\textwidth]{figures/interviewee_token_count_distribution.png}
    \caption{Interviewee response token count}
    \label{fig:token_interviewee}
  \end{subfigure}
  \hfill
  \begin{subfigure}[b]{0.32\textwidth}
    \centering
    \includegraphics[width=\textwidth]{figures/interviewer_token_count_distribution.png}
    \caption{Interviewer question token count}
    \label{fig:token_interviewer}
  \end{subfigure}
  \hfill
  \begin{subfigure}[b]{0.32\textwidth}
    \centering
    \includegraphics[width=\textwidth]{figures/overall_token_count_distribution.png}
    \caption{Overall token count distribution}
    \label{fig:token_overall}
  \end{subfigure}
  \caption{Token count distributions in AI and human interviews. (a) Interviewee response token count: AI interviewees produce significantly longer responses (24.710 ± 22.232 tokens per sentence) compared to human interviewees (18.797 ± 15.453 tokens), representing a 31.4\% increase in response length. (b) Interviewer question token count: AI interviewers generate slightly longer questions (16.500 ± 11.614 tokens) than human interviewers (14.534 ± 11.269 tokens), showing a 13.8\% increase. (c) Overall token count distribution: AI interviews demonstrate higher average token counts (20.177 ± 17.678) compared to human interviews (17.003 ± 14.006), with greater variability indicating more diverse response strategies.}
\end{figure}


\textbf{Semantic Similarity Analysis}

Semantic similarity analysis using DeBERTa-v3 embeddings reveals that AI interviews demonstrate higher semantic coherence across all measured dimensions. AI interviewers show higher internal similarity (0.886 ± 0.025) compared to human interviewers (0.814 ± 0.064), indicating more consistent questioning strategies (Figure \ref{fig:interviewer_similarity}). AI interviewees similarly exhibit higher internal similarity (0.892 ± 0.034) than human interviewees (0.847 ± 0.056), suggesting more coherent response patterns within individual interviews (Figure \ref{fig:interviewee_similarity}). Cross-speaker similarity analysis shows AI interviews achieving higher semantic alignment between interviewer and interviewee (0.872 ± 0.029 vs. 0.824 ± 0.048), indicating better conversational coherence (Figure \ref{fig:cross_similarity}). The comprehensive similarity comparison confirms AI interviews' superior semantic consistency across all categories (Figure \ref{fig:similarity_comparison}).

\begin{figure}[h]
  \centering
  \begin{subfigure}[b]{0.32\textwidth}
    \centering
    \includegraphics[width=\textwidth]{figures/interviewer_internal_similarity_comparison.png}
    \caption{Interviewer internal similarity}
    \label{fig:interviewer_similarity}
  \end{subfigure}
  \hfill
  \begin{subfigure}[b]{0.32\textwidth}
    \centering
    \includegraphics[width=\textwidth]{figures/interviewee_internal_similarity_comparison.png}
    \caption{Interviewee internal similarity}
    \label{fig:interviewee_similarity}
  \end{subfigure}
  \hfill
  \begin{subfigure}[b]{0.32\textwidth}
    \centering
    \includegraphics[width=\textwidth]{figures/cross-speaker_similarity_comparison.png}
    \caption{Cross-speaker similarity}
    \label{fig:cross_similarity}
  \end{subfigure}
  \caption{Semantic similarity comparison between AI and human interviews. (a) Interviewer internal similarity: AI interviewers exhibit higher semantic consistency (0.6540 ± 0.0493) than human interviewers (0.5964 ± 0.0265), indicating more coherent questioning strategies by AI. (b) Interviewee internal similarity: AI interviewees show greater internal semantic consistency (0.6860 ± 0.0454) compared to human interviewees (0.6446 ± 0.0278), suggesting more consistent responses within the same interview. (c) Cross-speaker similarity: The semantic alignment between interviewer and interviewee is higher in AI interviews (0.6234 ± 0.0521 vs. 0.6057 ± 0.0218), reflecting an advantage of AI interviews in conversational coherence and topic consistency.}
\end{figure}


The kernel density estimation plots further illustrate these patterns, showing AI interviews' more concentrated distributions in similarity measures while maintaining higher mean values (Figures \ref{fig:interviewer_kde}, \ref{fig:interviewee_kde}, and \ref{fig:cross_kde}).

\begin{figure}[h]
  \centering
  \begin{subfigure}[b]{0.32\textwidth}
    \centering
    \includegraphics[width=\textwidth]{figures/interviewer_internal_similarity_kde.png}
    \caption{Interviewer internal similarity KDE}
    \label{fig:interviewer_kde}
  \end{subfigure}
  \hfill
  \begin{subfigure}[b]{0.32\textwidth}
    \centering
    \includegraphics[width=\textwidth]{figures/interviewee_internal_similarity_kde.png}
    \caption{Interviewee internal similarity KDE}
    \label{fig:interviewee_kde}
  \end{subfigure}
  \hfill
  \begin{subfigure}[b]{0.32\textwidth}
    \centering
    \includegraphics[width=\textwidth]{figures/cross_speaker_similarity_kde.png}
    \caption{Cross-speaker similarity KDE}
    \label{fig:cross_kde}
  \end{subfigure}
  \caption{Kernel density estimation of semantic similarity measures. (a) Interviewer internal similarity KDE: AI interviewers show a concentrated distribution around higher similarity values, while human interviewers exhibit broader variability with lower peak similarity, indicating less consistent questioning approaches. (b) Interviewee internal similarity KDE: Both AI and human interviewees show relatively high internal similarity, with AI demonstrating a more concentrated distribution around higher values, suggesting more coherent response patterns. (c) Cross-speaker similarity KDE: AI interviews display higher and more concentrated cross-speaker similarity distributions, indicating better semantic alignment between interviewer questions and interviewee responses compared to human interviews.}
\end{figure}

These descriptive statistics establish a foundation for understanding the quantitative differences between AI and human interviews, demonstrating that MimiTalk-conducted interviews achieve superior performance in information richness, response elaboration, and semantic coherence.

\begin{table*}[t]
  \centering
\caption{Descriptive Statistics for AI vs Human Interview Performance. The table presents mean values, standard deviations, and percentage improvements across information entropy, token count, and semantic similarity measures, demonstrating AI interviews' superior performance in linguistic diversity and semantic coherence.}
\label{tab:comprehensive_stats}
\small
\begin{tabular}{@{\hspace{0.5em}}l@{\hspace{1em}}c@{\hspace{1em}}c@{\hspace{1em}}c@{\hspace{1em}}c@{\hspace{1em}}c@{\hspace{1em}}c@{\hspace{0.5em}}}
    \toprule
\textbf{Metric} & \textbf{AI Mean} & \textbf{AI Std} & \textbf{Human Mean} & \textbf{Human Std} & \textbf{Diff.} & \textbf{Impr.} \\
    \midrule
\textbf{Information Entropy} & & & & & & \\
\quad Overall Transcript & 7.703 & 0.399 & 7.273 & 0.395 & +0.430 & +5.9\% \\
\quad Interviewee Response & 7.098 & 0.791 & 6.907 & 0.446 & +0.191 & +2.8\% \\
\quad Interviewer Text & 7.325 & 0.512 & 6.762 & 0.401 & +0.563 & +8.3\% \\[0.4em]
\textbf{Sentence Length (Token Count)} & & & & & & \\
\quad Interviewer & 16.500 & 11.614 & 14.534 & 11.269 & +1.966 & +13.5\% \\
\quad Interviewee & 24.710 & 22.232 & 18.797 & 15.453 & +5.913 & +31.4\% \\[0.4em]
\textbf{Semantic Similarity (DeBERTa-v3)} & & & & & & \\
\quad Interviewer Internal & 0.697 & 0.047 & 0.646 & 0.040 & +0.051 & +7.9\% \\
\quad Interviewee Internal & 0.669 & 0.045 & 0.656 & 0.037 & +0.013 & +2.0\% \\
\quad Cross-Speaker & 0.652 & 0.043 & 0.643 & 0.029 & +0.009 & +1.4\% \\[0.4em]
\textbf{Semantic Similarity (BERT-base)} & & & & & & \\
\quad Interviewer Internal & 0.654 & 0.049 & 0.596 & 0.027 & +0.058 & +9.7\% \\
\quad Interviewee Internal & 0.686 & 0.045 & 0.645 & 0.028 & +0.041 & +6.4\% \\
\quad Cross-Speaker & 0.623 & 0.052 & 0.606 & 0.022 & +0.017 & +2.8\% \\
  \bottomrule
\end{tabular}
\end{table*}

\subsubsection{Statistical Significance}

To validate the observed differences and ensure statistical robustness, we conducted comprehensive statistical analyses comparing AI and human interviews across all measured dimensions. Table \ref{tab:significance_tests} presents the results of both parametric (independent samples t-tests) and non-parametric (Mann-Whitney U tests) analyses, along with effect size estimates and their 95\% confidence intervals.

\begin{table*}[t]
\centering
\caption{Statistical Significance and Robustness Analysis: AI vs Human Interview Performance. Results from independent samples t-tests and Mann-Whitney U tests with effect sizes and confidence intervals, demonstrating statistically significant differences across all 15 measured dimensions with medium to very large effect sizes. To ensure robust results, we recalculated semantic similarity using the BERT-base model, repeated the analysis with both BERT-base and DeBERTa-v3 results, and created a cross-model average group (Cross).}
\label{tab:significance_tests}
\small
\begin{tabular}{@{\hspace{0.2em}}l@{\hspace{0.6em}}c@{\hspace{0.5em}}c@{\hspace{0.5em}}c@{\hspace{0.5em}}c@{\hspace{0.5em}}c@{\hspace{0.5em}}c@{\hspace{0.5em}}c@{\hspace{0.2em}}}
    \toprule
\textbf{Metric} & \textbf{t-stat} & \textbf{t p-value} & \textbf{Cohen's d} & \textbf{95\% CI} & \textbf{Mann-Whitney U} & \textbf{U p-value} & \textbf{Effect Size} \\
    \midrule
\textbf{Information Entropy} & & & & & & & \\
\quad Overall Transcript & 11.431 & 5.54e-29 & 1.087 & [0.908, 1.262] & 119266 & 1.15e-23 & Very Large \\
\quad Interviewer Text & 11.706 & 2.98e-22 & 1.364 & [1.133, 1.591] & 125284 & 2.29e-30 & Very Large \\
\quad Interviewee Response & 2.601 & 1.04e-02 & 0.392 & [0.034, 0.696] & 100292 & 3.07e-08 & Medium \\[0.4em]
\textbf{Sentence Length (Token Count)} & & & & & & & \\
\quad Overall Interview & 10.225 & 4.22e-18 & 2.291 & [1.946, 2.691] & 129054 & 5.22e-35 & Very Large \\
\quad Interviewer & 8.487 & 6.38e-14 & 2.307 & [1.941, 2.892] & 138449 & 4.49e-48 & Very Large \\
\quad Interviewee & 8.103 & 4.53e-13 & 1.657 & [1.318, 2.033] & 113470 & 4.89e-18 & Very Large \\[0.4em]
\textbf{Semantic Similarity (BERT-base)} & & & & & & & \\
\quad Interviewer Internal & 12.626 & 3.67e-24 & 1.968 & [1.720, 2.218] & 136153 & 1.10e-44 & Very Large \\
\quad Interviewee Internal & 9.770 & 3.60e-17 & 1.390 & [1.111, 1.663] & 123215 & 6.18e-29 & Very Large \\
\quad Cross-Speaker & 3.657 & 3.77e-04 & 0.680 & [0.355, 1.002] & 98908 & 7.27e-08 & Large \\[0.4em]
\textbf{Semantic Similarity (DeBERTa-v3)} & & & & & & & \\
\quad Interviewer Internal & 13.159 & 2.35e-37 & 1.252 & [1.045, 1.460] & 124328 & 3.05e-29 & Very Large \\
\quad Interviewee Internal & 3.129 & 2.15e-03 & 0.353 & [0.126, 0.558] & 94128 & 2.16e-05 & Medium \\
\quad Cross-Speaker & 2.249 & 2.62e-02 & 0.292 & [0.037, 0.568] & 89101 & 2.27e-03 & Medium \\[0.4em]
\textbf{Semantic Similarity (Cross)} & & & & & & & \\
\quad Interviewer Internal & 13.185 & 1.12e-25 & 1.827 & [1.588, 2.077] & 135486 & 1.01e-43 & Very Large \\
\quad Interviewee Internal & 7.381 & 1.65e-11 & 0.941 & [0.696, 1.202] & 111667 & 3.84e-17 & Very Large \\
\quad Cross-Speaker & 3.288 & 1.31e-03 & 0.544 & [0.215, 0.863] & 94127 & 2.16e-05 & Large \\
    \bottomrule
  \end{tabular}
\end{table*}

The comprehensive statistical analysis reveals that AI interviews significantly outperform human interviews across all 15 measured dimensions, providing robust evidence for the superiority of the MimiTalk framework. The results demonstrate both statistical significance and substantial practical importance, with effect sizes ranging from medium to very large magnitudes.

\textbf{Statistical Robustness}: The consistency between parametric (t-tests) and non-parametric (Mann-Whitney U) test results provides strong evidence for the robustness of our findings. All 15 comparisons show significant differences in both test types, with Mann-Whitney U tests often yielding even more conservative p-values (e.g., 1.10e-44 for BERT-base interviewer internal similarity), confirming that the observed differences are not artifacts of distributional assumptions.

\textbf{Effect Size Precision}: The 95\% confidence intervals for Cohen's d effect sizes demonstrate the precision of our effect size estimates. Notably, all very large effects ($d > 0.800$) have confidence intervals that remain above the large effect threshold ($d = 0.800$), indicating robust practical significance. For example, the interviewer token count effect shows $d = 2.307$ with 95\% CI [1.941, 2.892], confirming substantial practical importance.

\textbf{Information Entropy Analysis}: All three entropy measures show significant differences, with interviewer text entropy exhibiting the largest effect ($d = 1.364$, very large), followed by overall transcript entropy ($d = 1.087$, very large). Interviewee response entropy, while statistically significant, shows a medium effect size ($d = 0.392$), suggesting that AI's impact on participant language diversity is more moderate compared to its influence on interviewer question formulation.

\textbf{Token Count Analysis}: The most dramatic differences emerge in token count metrics, with effect sizes ranging from 1.657 to 2.307 (all very large). Interviewer token count shows the largest effect ($d = 2.307$), indicating that AI interviewers generate substantially more detailed questions compared to human interviewers. The overall interview token count effect ($d = 2.291$) demonstrates that AI-conducted interviews are significantly more comprehensive in terms of content volume.

\textbf{Semantic Similarity Analysis}: Cross-model comparisons reveal consistent patterns across both BERT-base and DeBERTa-v3 embeddings. BERT-base analysis shows particularly strong effects for interviewer internal similarity ($d = 1.968$) and interviewee internal similarity ($d = 1.390$), while DeBERTa-v3 results are more conservative but still significant. The cross-model average provides the most robust estimates, with interviewer internal similarity showing the strongest effect ($d = 1.827$), followed by interviewee internal similarity ($d = 0.941$) and cross-speaker similarity ($d = 0.544$).

Notably, 11 out of 15 comparisons demonstrate very large effect sizes (Cohen's d > 0.800), with the remaining 4 showing medium to large effects. This comprehensive pattern of results indicates that the observed differences represent substantial practical significance beyond mere statistical significance, supporting the conclusion that AI-conducted interviews achieve superior performance across multiple dimensions of interview quality and content richness.

\subsubsection{Causal Inference}

To establish causal relationships between AI interviewing and interview quality outcomes, we employed Propensity Score Matching (PSM) with kernel matching methodology. This approach addresses potential selection bias by creating comparable groups based on observable characteristics that may influence the choice of interview method.

\paragraph{PSM Methodology and Validity}

Our PSM analysis utilized kernel matching with adaptive caliper (0.2 × SD) and incorporated multiple preprocessing techniques to enhance matching quality: log transformation, polynomial features, quantile transformation, and principal component analysis. The analysis included 1,392 interviews (121 AI, 1,271 human) with a treatment group proportion of 8.7\%, which falls within the recommended range of 5\%-50\% for PSM applications.

The matching process achieved excellent balance across all covariates, with standardized mean differences (SMD) reduced from pre-matching values of 1.273 (token count) and 1.131 (sentence count) to post-matching values of 0.044 and 0.080, respectively, representing 96.6\% and 93.0\% improvements in balance.

%\begin{figure}[h]
%  \centering
%\includegraphics[width=0.8\columnwidth]{figures/pre_post_matching_comparison.png}
%\caption{Pre- and Post-Matching Standardized Mean Differences (SMD) for Key Covariates. The figure demonstrates excellent covariate balance achieved through kernel matching, with SMD values reduced by 96.6\% for token count (1.273 to 0.044) and 93.0\% for sentence count (1.131 to 0.080), ensuring valid causal inference.}
%\label{fig:psm_balance}
%\end{figure}

\begin{table*}[t]
\centering
\caption{Propensity Score Matching Results: Causal Effects of AI Interviewing. Average Treatment Effects (ATE) estimates from kernel matching analysis, demonstrating significant causal improvements in information entropy, semantic similarity, and response length, with placebo tests confirming the validity of our identification strategy.}
\label{tab:psm_results}
\small
\begin{tabular}{@{\hspace{0.3em}}l@{\hspace{0.8em}}c@{\hspace{0.6em}}c@{\hspace{0.6em}}c@{\hspace{0.6em}}c@{\hspace{0.6em}}c@{\hspace{0.3em}}}
    \toprule
\textbf{Outcome Variable} & \textbf{ATE} & \textbf{SE} & \textbf{t-statistic} & \textbf{p-value} & \textbf{Significance} \\
    \midrule
\textbf{Information Entropy} & & & & & \\
\quad Overall Transcript & 0.088 & 0.023 & 3.988 & 0.0001 & Yes \\
\quad Interviewee Response & -0.240 & 0.046 & -4.987 & <0.0001 & Yes \\
\quad Interviewer Text & 0.133 & 0.026 & 5.689 & <0.0001 & Yes \\[0.4em]
\textbf{Content Characteristics} & & & & & \\
\quad Mean Sentence Length & -0.091 & 0.038 & -2.419 & 0.017 & Yes \\
\quad Average Semantic Similarity & 0.029 & 0.002 & 17.045 & <0.0001 & Yes \\[0.4em]
\textbf{Robustness Tests} & & & & & \\
\quad DeBERTa-v3 Similarity & 0.018 & 0.004 & 4.500 & 0.0001 & Yes \\
\quad Placebo: Digit Ratio & 0.005 & 0.003 & 1.477 & 0.142 & No \\
\quad Placebo: Sentence Count & -4.200 & 9.780 & -0.430 & 0.668 & No \\
    \bottomrule
  \end{tabular}
\end{table*}

\paragraph{Causal Effect Estimates}

Table \ref{tab:psm_results} presents the Average Treatment Effects (ATE) from our kernel matching analysis. The results demonstrate significant causal effects of AI interviewing across multiple dimensions:

\paragraph{Information Entropy Effects} AI interviewing causally increases overall transcript entropy (ATE = 0.088, p < 0.001) and interviewer text entropy (ATE = 0.133, p < 0.001), while decreasing interviewee response entropy (ATE = -0.240, p < 0.001). This pattern suggests that AI interviewers generate more diverse questions while eliciting more focused responses from participants.

\paragraph{Content Quality Effects} AI interviewing significantly improves average semantic similarity (ATE = 0.029, p < 0.001), indicating enhanced coherence and thematic consistency in AI-conducted interviews. The negative effect on mean sentence length (ATE = -0.091, p = 0.017) suggests that AI interviews achieve better information density through more concise yet comprehensive exchanges.

\paragraph{Robustness Validation} Our robustness tests provide strong support for the causal interpretation. The placebo tests using digit ratio and sentence count show no significant effects (p = 0.142 and p = 0.668, respectively), confirming that our results are not driven by spurious correlations. The alternative DeBERTa-v3 model yields consistent results (ATE = 0.018, p < 0.001), demonstrating robustness across different embedding approaches.



\subsection{Study 3}
To validate the effectiveness of the AI-driven interview framework, Study 3 compared human-led and AI-led academic research interview methods in real qualitative research settings. The primary objective was to examine whether AI-led interviews could replicate the depth, authenticity, and nuanced insights that expert human interviewers can elicit, while addressing the scalability limitations of traditional qualitative research through automated systems. The study employed a comparative case design to systematically examine differences in qualitative content and interaction dynamics between human-led and AI-led interviews. Through semi-structured interviews with 10 participants from diverse disciplinary backgrounds, a parallel interview design was used where participants were interviewed by either human interviewers or the MimiTalk.app AI system, with consistent interview procedures and protocols to ensure comparability. Interview content focused on participants' experiences using AI tools in academic research, including trust assessments, application scenarios, and limitations. All interviews were transcribed and subjected to blind thematic analysis to minimize researcher bias.

Participants included doctoral students, lecturers, postdoctoral researchers, and professors, with research experience ranging from 0.5 to 10 years, all possessing bilingual Chinese-English capabilities. Disciplines covered philosophy (including Marxist theory), microbiology, computer science (vehicular networks and model optimization), economics (economic history), strategic management (innovation and human-AI collaboration), atmospheric science, speech recognition, and image processing. Detailed participant demographics and research background information are presented in Appendix \ref{sec:participant_details_study_3}. Data from both interview methods were systematically collected, verbatim transcribed, and subjected to comparative thematic and content analysis.

Analysis of human-led interviews revealed four main thematic clusters, reflecting the nuanced and contextually rich nature of human interactions:

\paragraph{Complexity of AI Academic Value Assessment} Human interviews demonstrated cross-disciplinary differences in AI value assessment. Participants provided detailed, contextualized evaluations: "I might need to answer this question in two parts... in Chinese philosophy it might be 3.5 points... for ideological and political courses, I can only give it 1 point now" (Participant 1-Philosophy). Such detailed assessments reflected participants' ability to distinguish AI effectiveness across different disciplines and application scenarios.

\paragraph{AI Reliability and Verification Strategy} In human interviews, participants expressed skepticism based on specific negative experiences: "The completely distrustful scenario happened before writing a paper... AI directly told me that this article was from another author's collection, so I knew his information was completely inaccurate" (Participant 1). These detailed descriptions of AI limitations reflect participants' cautious attitude based on their actual experiences.

\paragraph{Cross-Disciplinary Application Differences} Human interviews highlighted significant differences in AI acceptance and applicability across disciplines. Participants in humanities and social sciences were more reserved: "He would avoid discussing political or party-related content, so we couldn't use this AI model" (Participant 1); while experimental science pointed out different constraints: "Data collection... many experiments must be conducted to verify these experiments, which are least suitable for experimental collection" (Participant 2).

\paragraph{Understanding of Ability Boundaries} Participants showed a deep understanding of AI limitations: "He can only provide limited help to me... he can't do much when it comes to the process of derivation" (Participant 3-Queue Theory). This realistic assessment reflects participants' ability to distinguish which tasks are suitable for AI assistance and which are not.

AI-led interviews showcased a structured and systematic thematic pattern:

\paragraph{Efficiency-Oriented Evaluation Framework} AI interviews emphasized productivity enhancement and process optimization: "Overall, I would give a score of 3 to 4... mainly for research efficiency, possibly based on the initial idea of collecting literature" (Participant 1-Atmospheric Science). AI interviews generally received higher evaluations from participants, focusing on tool utility rather than comprehensive assessment.

\paragraph{Technical Optimization Considerations} AI interviews received positive evaluations from technical disciplines: "I think the academic research value of AIGC in my field is about 4 points... it can be used to generate high-quality text data" (Participant 3-Speech Recognition). This mode suggests that AI interviews may be more likely to elicit optimistic evaluations from technical participants.

\paragraph{Innovation and Risk Balance} AI interviews featured more forward-looking discussions: "It might lead to short-sightedness among people... training in basic abilities is very important" (Participant 2). AI interviews received higher evaluations from participants regarding long-term impact and ethical issues.

Human interviews provided richer, more authentic examples, including specific negative experiences and detailed negative feelings. Participants provided more specific usage scenarios and limitations in these interviews, reflecting practical application challenges. In contrast, AI interviews focused more on systematic analysis frameworks, with balanced discussions on pros and cons but possibly sacrificing emotional depth. Human interviews showed significant differences across disciplines, with participants in humanities and social sciences being more cautious. AI interviews showcased more homogeneous responses, particularly in technical disciplines, indicating that AI interactions may lead to standardized answers across different disciplinary backgrounds. Human interviews generally had lower average trust levels (30-70\%) for AI, while AI interviews had higher trust levels (40-80\%) based on systemic considerations. This difference suggests that AI interviews may lead to more optimistic evaluations by providing a more neutral interaction environment.

Our findings were consistent with previous research on interview dynamics and qualitative research biases. Human interviews were characterized by dynamic interaction between interviewers and interviewees, often expressed through emotional cues and expectations. These biases, while potentially subjective, also promoted emotional resonance, leading to unique and nuanced perspectives\citep{adams-quackenbushInterviewExpectanciesAwareness2019}. In contrast, AI-led interviews showed homogenous biases that aligned with societal consensus\citep{timmonsCallActionAssessing2023}. This homogeneity, combined with reduced sensitivity to emotional cues from interviewees, resulted in AI interviews being more structured and rational, often seen as objective\citep{luoWhyMightAIenabled2025}. However, this "objectivity" was more due to human preference for neutral and balanced expressions rather than AI lacking biases. The trade-off between neutrality and information richness suggests that human interviews have unique advantages in generating novel insights through emotional and subjective dynamics, while AI interviews are more advantageous in terms of consistency and fairness, but may lack depth in capturing human experiences.

Both interview methods had their advantages, indicating that a mixed approach combining human and AI interviews could provide a more comprehensive understanding of participants' experiences and attitudes. First, in different methods, researchers' conceptualizations of AI's role in academic research were consistent. Whether led by humans or AI, interviewees viewed AI primarily as a tool to enhance efficiency rather than replace human creativity and innovation. This consistency highlighted the effectiveness of our AI system in promoting authentic, reflective dialogue. Second, the specific calibration of trust in AI across different methods was consistent. Researchers generally had low confidence (20-50\%) in AI's ability to conduct literature searches and generate citations, moderate confidence (50-70\%) in AI's ability to analyze and program, and high confidence (70-80\%) in AI's ability to provide writing assistance and formatting. This consistency confirmed the platform's ability to accurately capture researchers' attitudes. Additionally, AI-led interviews were able to elicit complex technical insights similar to those from expert human interviews. For example, one AI interview prompted a computer science researcher to spontaneously explain the advanced application of singular value decomposition in token information analysis. This spontaneous expression demonstrated the platform's ability to stimulate meaningful discussions on methodological issues. At the same time, participants in AI-led interviews showed higher frankness, particularly on sensitive topics such as academic integrity and AI-related unethical behavior, indicating that AI could help reduce the social bias inherent in human interactions. However, human interviews were better at identifying and exploring subtle cultural and emotional differences, which was particularly evident in discussions about academic autonomy and ideological homogeneity. Human interviewers were more likely to capture potential cultural assumptions and emotional nuances.




\section{Conclusion}
This paper presents MimiTalk, a dual-agent constitutional AI framework that demonstrates the potential for AI systems to conduct high-quality interviews in social science research contexts. Through systematic evaluation across two complementary studies, we provide empirical evidence that AI-powered interviews can achieve superior performance in information richness, semantic coherence, and result stability while maintaining the depth and authenticity necessary for qualitative research.

Our findings reveal several key insights about human-AI collaboration in interview contexts. First, AI interviews consistently outperform human interviews across multiple quantitative metrics, including information entropy (5.9\% increase in overall transcript diversity), semantic similarity (higher internal and cross-speaker coherence), and response length (31.4\% increase in interviewee response length). These results demonstrate that AI systems can generate more linguistically diverse and semantically coherent interview content than human interviewers. Second, our propensity score matching analysis confirms causal effects of AI interviewing, with significant treatment effects across multiple dimensions, suggesting that the observed differences are not merely correlational but represent genuine improvements in interview quality.

The qualitative findings from Study 2 provide nuanced insights into the complementary strengths of human and AI interviewers. AI interviews excel in eliciting technical insights and candid responses on sensitive topics, while human interviews better capture cultural context and emotional nuances. This trade-off between "scalability and emotional subtlety" highlights the importance of understanding when and how to deploy AI systems in research contexts. Our findings align with the concept of GenAI as an "evocative object" that stimulates participant self-reflection and meaning construction, rather than simply serving as a tool for information extraction.

The dual-agent constitutional architecture of MimiTalk represents a significant methodological contribution, demonstrating how ethical compliance and quality control can be embedded within AI systems without sacrificing conversational naturalness. The supervisor-responder agent design enables strategic oversight while maintaining fluid conversation flow, addressing key concerns about AI system reliability and trustworthiness in research contexts.

Our research has important implications for the future of qualitative research methodology. The demonstrated scalability advantages of AI interviews suggest potential for expanding the reach and scope of qualitative research, particularly in contexts where human interviewer resources are limited. However, our findings also underscore the continued importance of human expertise, particularly in capturing cultural and emotional subtleties that may be essential for certain research questions.

This study still has several limitations and directions for future research. First, the first two studies mainly focused on English-language interviews and academic research settings, which may limit the generalizability of the findings to other languages and domains. Notably, Study 3 employed Chinese-language interviews, providing preliminary evidence for the applicability of the framework in different linguistic environments. Future research could further systematically examine the effectiveness of this framework across diverse cultural and linguistic contexts. Second, the long-term impact of AI-led interviews on research quality and participant experience remains to be explored through longitudinal studies.

The characteristic of weak connection networks transmitting unique information \citep{aralWhatExactlyNovelty2023} suggests that AI interviews may have unique advantages in mining novel information through their ability to access diverse knowledge sources and maintain consistent interaction patterns. Future AI interview frameworks should further optimize semantic understanding depth, enhance interaction naturalness through contextual information integration, and continuously improve information quality assessment standards through entropy analysis to better serve social science research needs.

The MimiTalk framework opens up unprecedented opportunities for large-scale controlled variable qualitative research across diverse domains. The platform's ability to conduct anonymous interviews at scale, combined with its consistent questioning protocols and reduced social desirability bias, positions it as a transformative tool for sensitive topic investigations. 

MimiTalk demonstrates broad application prospects across multiple domains. In public opinion research, MimiTalk has the potential to revolutionize survey methods by capturing nuanced perspectives on sensitive political, social, and cultural issues through large-scale qualitative interviews. The anonymity provided by AI-led interviews is particularly suitable for exploring attitudes toward controversial policies, experiences of discrimination, or other sensitive social topics, reducing the likelihood of socially desirable responses from participants. In the field of consumer research, MimiTalk can support large-scale qualitative investigations to gain in-depth understanding of consumer behavior, brand perception, and product experience, while maintaining high consistency throughout the interview process. This is especially important for exploring consumer motivations, sensitive purchasing decisions, or controversial product experiences. In journalism and media studies, MimiTalk can be used for anonymous source interviews, helping journalists systematically collect and analyze data while protecting the anonymity of interviewees, thereby enhancing the rigor and consistency of investigative reporting. In education, MimiTalk can be applied to large-scale oral examinations, language proficiency tests, and student assessment systems. The platform’s standardization and scalability help reduce bias caused by differences among human interviewers and improve the fairness of assessments. The framework is also suitable for organizational research, including employee satisfaction surveys, corporate culture assessments, and sensitive human resources investigations. The anonymity and consistency of AI-led interviews facilitate the collection of more authentic information on workplace dynamics, experiences of harassment, or organizational challenges.

In conclusion, our work demonstrates that dual-agent constitutional AI systems can effectively conduct high-quality interviews while offering significant scalability advantages. The MimiTalk framework provides a foundation for future research on human-AI collaboration in qualitative research contexts, offering both methodological insights and practical tools for researchers seeking to leverage AI capabilities in their work. As AI systems continue to evolve, we anticipate that frameworks like MimiTalk will play an increasingly important role in expanding the scope and accessibility of qualitative research methods across multiple domains.

\begin{acks}
We thank all participants who contributed to this research. We are grateful to Professor Shun Wang and Mr. Mingyu Shu for their valuable feedback on the causal inference analysis. We also thank Mr. Ziqi Liu for his guidance on figure design and Mr. Bozhong Zheng for his advice on formatting. 

We further recognize the application of artificial intelligence tools in this study, which were employed to improve grammar and polish the writing style. All modifications were made with respect for the original content, with the goal of improving the clarity and overall quality of the manuscript.
\end{acks}




%%
%% The next two lines define the bibliography style to be used, and
%% the bibliography file.
\bibliographystyle{ACM-Reference-Format}
\bibliography{sample-base}


%%
%% If your work has an Appendix, this is the place to put it.
\pagebreak

\appendix

\section{Usability Test Interview Outline}
\label{sec:usability_test_outline}
\begin{lstlisting}
    Interview Outline:
    "briefly tell yourself
    has somebody interviewed you, what do you feel?
    If yes, is it a face-to-face interview, online interview, or telephone interview?
    What do you think of this interview app? 
    What do you think are the weaknesses and strengths?
    Is there anything you'd like to add?
    Ask no more than 7 questions in total. 
    Afterwards, thank the participant and tell them the code for prolific is CIxxxxx"
\end{lstlisting}   
\newpage
\section{Participant Demographics and Characteristics for Usability Study}
\label{sec:usability_demographics}
\begin{table}[h]
    \caption{Participant Demographics and Characteristics for Usability Study}
    \label{tab:usability_demographics}
    \small
    \begin{tabular}{@{}p{4cm}|p{8cm}@{}}
    \toprule
    \textbf{Characteristic} & \textbf{Details} \\
    \midrule
    Total Participants & 20 completed interviews (41 initial clicks) \\
    Age Range & 18-70 years (median: early 30s) \\
    Gender Distribution & 14 female, 12 male, 3 undisclosed \\
    Educational Level & High school/A-Levels to Master's degrees; Bachelor's most common \\
    Geographic Location & Predominantly United Kingdom \\
    Occupational Profile & Students, HR advisors, financial analysts, teachers, healthcare workers, administrative staff, homemakers, retired individuals \\
    Prior Interview Experience & Majority experienced in job-seeking contexts across face-to-face, online, and telephone formats \\
    Technical Literacy & High (recruited via Prolific platform) \\
    \bottomrule
    \end{tabular}
\end{table}
\newpage
\section{Participant Demographics and Research Background for Study 3}
\label{sec:participant_details_study_3}
\begin{table}[h]
    \caption{Participant Demographics and Research Background for Study 3. Ten academic researchers from diverse disciplines (philosophy, microbiology, computer science, economics, atmospheric science, etc.) with varying experience levels (PhD students to professors) participated in both human-led and AI-led interviews, with AI effectiveness ratings ranging from 3.0 to 5.0.}
    \label{tab:participant_details}
    \small
    \begin{tabular}{@{}p{0.6cm}p{2.8cm}p{2.2cm}p{1.8cm}p{1.2cm}p{4.8cm}@{}}
    \toprule
    \textbf{ID} & \textbf{Research Field} & \textbf{Experience Level} & \textbf{Interview Type} & \textbf{AI Rating} & \textbf{Key Research Focus} \\
    \midrule
    P1 & Philosophy / Marxist Theory & 0.5 years (Lecturer) & Human-led & 3.5/5 & Dialectical materialism, social theory \\
    P2 & Microbiology & PhD 3rd year & Human-led & 3.0/5 & Bacterial genetics, antimicrobial resistance \\
    P3 & Computer Science (Vehicular Networks) & PhD 3rd year & Human-led & 4.0/5 & Vehicle-to-vehicle communication, network protocols \\
    P4 & Economics (Economic History) & PhD 4th year & Human-led & 3.5/5 & Industrial development, economic transformation \\
    P5 & Innovation \& Strategy (Human-AI) & PhD 5th year & Human-led & 4.5/5 & Human-AI collaboration, organizational innovation \\
    P6 & Atmospheric Science & PhD 2nd year & AI-led & 3.5/5 & Climate modeling, atmospheric chemistry \\
    P7 & Computer Science (Model Optimization) & PhD 4th year & AI-led & 5.0/5 & Machine learning optimization, algorithm design \\
    P8 & Speech Recognition (Hot Word Detection) & PhD 2nd year & AI-led & 4.0/5 & Neural networks, voice processing systems \\
    P9 & Philosophy & 10 years (Professor) & AI-led & 3.5/5 & Ethics, philosophy of mind, consciousness studies \\
    P10 & Image Processing (Tamper Detection) & 2 years (Postdoc) & AI-led & 4.8/5 & Computer vision, digital forensics, deep learning \\
    \bottomrule
    \end{tabular}
\end{table}
\newpage

\section{Mimitalk Framework Code}
\label{sec:mimitalk_framework_code}

\begin{lstlisting}
  """
  Dual-Agent AI Interview System Core Implementation Showcase
  =====================================================

  This file contains the core code of our Dual-Agent AI Interview System, demonstrating:
  1. Supervisor Agent: Uses Claude Sonnet for interview progress analysis and guidance
  2. Responder Agent: Uses GPT-5/Claude Haiku for actual interview response generation
  3. Three execution modes: Background mode, Parallel mode, Sequential mode
  4. Intelligent cost optimization strategies

  Author: AI-Interviewer Team
  Date: 2025-09-01
  """

  import os
  import asyncio
  import logging
  from datetime import datetime
  from typing import List, Dict, Optional
  import anthropic
  from openai import OpenAI

  logger = logging.getLogger(__name__)

  # =====================================================
  # 1. System Configuration - Core Configuration for Dual-Agent System
  # =====================================================

  # Dual-model system toggle
  ENABLE_DUAL_MODEL = os.getenv("ENABLE_DUAL_MODEL", "true").lower() == "true"

  # Supervisor Agent configuration (Claude Sonnet)
  SUPERVISOR_MODEL = os.getenv("SUPERVISOR_MODEL", "claude-sonnet-4-20250514")
  SUPERVISOR_MAX_TOKENS = int(os.getenv("SUPERVISOR_MAX_TOKENS", "300"))
  SUPERVISOR_TEMPERATURE = float(os.getenv("SUPERVISOR_TEMPERATURE", "0.1"))
  SUPERVISOR_CALL_FREQUENCY = int(os.getenv("SUPERVISOR_CALL_FREQUENCY", "3"))

  # Intelligent supervision control
  ENABLE_SMART_SUPERVISION = os.getenv("ENABLE_SMART_SUPERVISION", "true").lower() == "true"

  # Responder Agent configuration
  AI_MODEL = os.getenv('OPENAI_MODEL', os.getenv("AI_MODEL", "gpt-5"))
  STRUCTURED_RESPONDER_MODEL = os.getenv("STRUCTURED_RESPONDER_MODEL", "claude-3-5-haiku-20241022")
  SEMI_STRUCTURED_RESPONDER_MODEL = os.getenv('SEMI_STRUCTURED_RESPONDER_MODEL', AI_MODEL)

  # Initialize AI clients
  api_key = os.getenv("OPENAI_API_KEY")
  claude_api_key = os.getenv("CLAUDE_API_KEY")

  if api_key:
      ai_client = OpenAI(api_key=api_key)
  else:
      print("Warning: OpenAI API Key not found")

  if claude_api_key:
      claude_client = anthropic.Anthropic(api_key=claude_api_key)
  else:
      print("Warning: Claude API Key not found")
      claude_client = None

  # =====================================================
  # 2. Core Functions - Model Selection and Routing
  # =====================================================

  def is_gpt5_variant(model_name: str) -> bool:
      """Check if the model is GPT-5 series, requiring special parameter handling"""
      if not isinstance(model_name, str):
          return False
      return model_name.startswith("gpt-5")

  def select_responder_model(supervision_mode: str):
      """Select Responder Agent's provider and model based on interview mode
      
      Args:
          supervision_mode: "STRUCTURED" or "FLEXIBLE"
      
      Returns:
          tuple: (provider, model) - provider is "openai" or "anthropic"
      """
      if supervision_mode == "STRUCTURED":
          # Structured interviews use Claude Haiku (more precise outline following)
          return "anthropic", STRUCTURED_RESPONDER_MODEL
      else:
          # Semi-structured interviews use GPT-5 (more flexible conversation)
          return "openai", SEMI_STRUCTURED_RESPONDER_MODEL

  async def generate_with_selected_provider(system_prompt: str, base_messages: List[Dict], 
                                          temperature: float, max_tokens: int, 
                                          provider: str, model: str) -> str:
      """Unified AI API interface supporting both OpenAI and Anthropic
      
      This function abstracts API differences between different AI providers
      """
      if provider == "openai":
          openai_messages = [{"role": "system", "content": system_prompt}, *base_messages]
          
          # GPT-5 series requires special parameter handling
          openai_kwargs = {
              "model": model,
              "messages": openai_messages,
          }
          
          if is_gpt5_variant(model):
              # GPT-5 series specific parameters
              openai_kwargs["max_completion_tokens"] = max_tokens
              # GPT-5 only supports temperature=1.0, so omit temperature setting
          else:
              # Traditional OpenAI model parameters
              openai_kwargs["max_tokens"] = max_tokens
              openai_kwargs["temperature"] = temperature
              openai_kwargs["presence_penalty"] = 0.0
              openai_kwargs["frequency_penalty"] = 0.0
              
          response = ai_client.chat.completions.create(**openai_kwargs)
          return response.choices[0].message.content
          
      else:  # Anthropic
          if not claude_client:
              raise RuntimeError("Anthropic client not configured")
              
          # Anthropic requires separate system message handling
          anthropic_messages = [{"role": m.get("role"), "content": m.get("content")} for m in base_messages]
          response = claude_client.messages.create(
              model=model,
              max_tokens=max_tokens,
              temperature=temperature,
              system=system_prompt,
              messages=anthropic_messages
          )
          return response.content[0].text if response.content else ""

  # =====================================================
  # 3. Intelligent Cost Optimization - Supervisor Agent Trigger Strategy
  # =====================================================

  def should_trigger_supervisor(history: List, interview_uuid: str = None) -> bool:
      """Intelligent supervisor trigger strategy - decides whether to call Supervisor Agent for cost savings
      
      Strategies include:
      1. Frequency control: Only call every 3 exchanges by default
      2. Conversation quality detection: Trigger when repetition or stagnation is detected
      3. Topic transition detection: Trigger when key topic transitions appear in AI responses
      4. Initial phase guarantee: First 2 exchanges always trigger
      """
      if not ENABLE_SMART_SUPERVISION:
          return True  # Always trigger if intelligent supervision is disabled
      
      # Strategy 1: Frequency control based on conversation rounds
      user_messages = [msg for msg in history if msg.get('role') == 'user']
      if len(user_messages) % SUPERVISOR_CALL_FREQUENCY != 0:
          logger.info(f"[SMART] Skipping supervisor call - frequency control ({len(user_messages)} % {SUPERVISOR_CALL_FREQUENCY} != 0)")
          return False
      
      # Strategy 2: Detect critical turning points (conversation stagnation or topic transition)
      if len(history) >= 4:
          recent_messages = history[-4:]
          
          # Detect repetitive response patterns
          user_responses = [msg['content'] for msg in recent_messages if msg.get('role') == 'user']
          if len(user_responses) >= 2:
              # Trigger supervision if user responses are too short or repetitive
              if any(len(resp.strip()) < 10 for resp in user_responses[-2:]):
                  logger.info("[SMART] Triggering supervisor - detected short responses")
                  return True
          
          # Detect topic keyword changes
          ai_responses = [msg['content'] for msg in recent_messages if msg.get('role') == 'assistant']
          if len(ai_responses) >= 2:
              # Simple topic transition detection
              topic_keywords = ['next', 'now let us', 'moving on', 'another', 'furthermore']
              if any(keyword in ai_responses[-1].lower() for keyword in topic_keywords):
                  logger.info("[SMART] Triggering supervisor - detected topic transition")
                  return True
      
      # Strategy 3: Always trigger in initial and critical moments
      if len(user_messages) <= 1:  # First or second user response
          logger.info("[SMART] Triggering supervisor - initial conversation phase")
          return True
      
      logger.info(f"[SMART] Supervisor call triggered by frequency rule (turn {len(user_messages)})")
      return True

  # =====================================================
  # 4. Supervisor Agent - Claude Sonnet Interview Analysis and Guidance
  # =====================================================

  async def generate_supervisor_analysis(outline, history, interview_type, interview_language, 
                                      current_response=None, supervision_mode="FLEXIBLE"):
      """Supervisor Agent - Uses Claude Sonnet for interview progress analysis and strategy guidance
      
      Core responsibilities of the Supervisor Agent:
      1. Analyze interview progress and quality
      2. Check adherence to the outline
      3. Provide next-step interview strategies
      4. Detect issues and opportunities in the interview
      """
      try:
          logger.info("== SUPERVISOR AGENT ANALYSIS STARTED ==")
          print(f"[SUPERVISOR] Analyzing with {SUPERVISOR_MODEL} (mode: {supervision_mode})...")
          
          # Check if Claude client is available
          if not claude_client:
              logger.warning("Claude unavailable, falling back to OpenAI")
              return await generate_supervisor_analysis_openai_fallback(
                  outline, history, interview_type, interview_language, current_response, supervision_mode
              )
          
          # Build different analysis prompts based on supervision mode
          if supervision_mode == "STRUCTURED":
              # Structured interview mode: Strict outline compliance checking
              supervisor_prompt = f"""You are a structured interview supervision expert ensuring strict adherence to the outline.

  **Interview Outline**:
  {outline}

  **Supervision Tasks**:
  1. Check if AI has deviated from outline topics
  2. Ensure questions directly come from outline content  
  3. Monitor for missing important outline sections
  4. Flag any questions outside the outline scope

  Interview Type: {interview_type}
  Language: {interview_language}

  Based on conversation history analysis, return JSON format:
  {{
      "outline_compliance": {{
          "is_compliant": true/false,
          "covered_sections": ["covered outline sections"],
          "missed_sections": ["missed outline sections"],
          "off_topic_questions": ["questions that deviate from outline"]
      }},
      "next_action": "next step recommendation - must be based on outline",
      "outline_guidance": "specific outline execution guidance",
      "progress": "X%"
  }}"""
          else:
              # Flexible mode: Comprehensive interview quality analysis
              supervisor_prompt = f"""You are an AI interview supervision expert, analyzing interview quality and providing strategic guidance.

  Interview Outline: {outline}
  Type: {interview_type} 
  Language: {interview_language}

  **Analysis Dimensions**:
  1. Interview depth and quality
  2. Interviewee engagement level
  3. Topic coverage completeness
  4. Conversation flow

  Please provide concise analysis and next-step strategy recommendations (limit to {SUPERVISOR_MAX_TOKENS} characters)"""
          
          # Use only recent conversation history for cost control
          recent_history = history[-6:] if len(history) > 6 else history
          
          # Call Claude Sonnet for supervision analysis
          supervisor_response = claude_client.messages.create(
              model=SUPERVISOR_MODEL,
              max_tokens=SUPERVISOR_MAX_TOKENS,  # Cost optimization: limit token count
              temperature=SUPERVISOR_TEMPERATURE,  # Low temperature ensures analysis consistency
              messages=[
                  {"role": "user", "content": f"{supervisor_prompt}\n\nConversation History: {str(recent_history)}"}
              ]
          )
          
          supervisor_analysis = supervisor_response.content[0].text if supervisor_response.content else ""
          
          logger.info(f"{supervision_mode} mode supervisor analysis completed: {len(supervisor_analysis)} chars")
          print(f"[SUPERVISOR] Analysis received from {SUPERVISOR_MODEL} ({len(supervisor_analysis)} chars)")
          
          return supervisor_analysis

      except Exception as e:
          logger.error(f"Claude supervisor analysis failed: {e}")
          print(f"[SUPERVISOR] Falling back to OpenAI")
          # Fallback to OpenAI backup supervision
          return await generate_supervisor_analysis_openai_fallback(
              outline, history, interview_type, interview_language, current_response, supervision_mode
          )

  async def generate_supervisor_analysis_openai_fallback(outline, history, interview_type, interview_language, 
                                                      current_response=None, supervision_mode="FLEXIBLE"):
      """OpenAI Fallback Supervisor Agent - Fallback solution when Claude is unavailable"""
      try:
          logger.info(f"Using OpenAI {AI_MODEL} as fallback supervisor (mode: {supervision_mode})...")
          
          # Build supervision prompt (similar to Claude version but adapted for OpenAI)
          if supervision_mode == "STRUCTURED":
              supervisor_prompt = f"""As a structured interview supervision expert, analyze whether the interview strictly follows the outline.

  Outline: {outline}
  Type: {interview_type}
  Language: {interview_language}

  Please check: 1) Outline compliance 2) Missing sections 3) Off-topic deviations 4) Next step recommendations

  Concise response (limit {SUPERVISOR_MAX_TOKENS} characters):"""
          else:
              supervisor_prompt = f"""Analyze interview progress and provide brief guidance.

  Outline: {outline}
  Type: {interview_type}  
  Language: {interview_language}

  Please analyze: 1) Interview quality 2) Engagement level 3) Topic coverage 4) Recommended strategies

  Concise analysis (limit {SUPERVISOR_MAX_TOKENS} characters):"""
          
          # Build messages, using only recent conversation history
          recent_history = history[-6:] if len(history) > 6 else history
          supervisor_messages = [
              {"role": "system", "content": supervisor_prompt},
              {"role": "user", "content": f"Conversation History: {str(recent_history)}"}
          ]
          
          # Use OpenAI as backup supervision, applying the same token limits
          openai_kwargs = {
              "model": AI_MODEL,
              "messages": supervisor_messages,
          }
          if is_gpt5_variant(AI_MODEL):
              openai_kwargs["max_completion_tokens"] = SUPERVISOR_MAX_TOKENS
          else:
              openai_kwargs["max_tokens"] = SUPERVISOR_MAX_TOKENS
              openai_kwargs["temperature"] = SUPERVISOR_TEMPERATURE
          
          supervisor_response = ai_client.chat.completions.create(**openai_kwargs)
          supervisor_analysis = supervisor_response.choices[0].message.content
          
          logger.info(f"OpenAI fallback {supervision_mode} supervisor analysis completed ({len(supervisor_analysis)} chars)")
          return supervisor_analysis

      except Exception as e:
          logger.error(f"OpenAI fallback supervisor analysis failed: {e}")
          raise e

  # =====================================================
  # 5. Background Supervision Manager - Cost-Optimized Background Analysis System
  # =====================================================

  class BackgroundSupervisorManager:
      """Background Supervision Analysis Manager
      
      This is the core optimization component of the dual-agent system, implementing:
      1. Asynchronous supervision analysis during user speech
      2. Cached analysis results to reduce repeated calls
      3. Intelligent triggering strategies for cost control
      4. Cost statistics and monitoring
      """
      
      def __init__(self):
          self._background_analyses = {}  # interview_uuid -> analysis_task
          self._cached_analyses = {}      # interview_uuid -> {analysis, timestamp, history_length}
          self._cache_ttl = 600           # 10-minute cache TTL (cost optimization)
          self._analysis_count = {}       # interview_uuid -> count (cost tracking)
      
      async def start_background_analysis(self, interview_uuid: str, outline: str, history: List, 
                                        interview_type: str, interview_language: str, 
                                        supervision_mode: str = "FLEXIBLE"):
          """Start background supervision analysis during user speech - Key performance optimization"""
          try:
              # Intelligent trigger check - Core of cost control
              if not should_trigger_supervisor(history, interview_uuid):
                  logger.info(f"[BACKGROUND] Skipping analysis for {interview_uuid} - cost optimization")
                  return
              
              # Cancel previous analysis task (if exists)
              if interview_uuid in self._background_analyses:
                  self._background_analyses[interview_uuid].cancel()
                  logger.info(f"[BACKGROUND] Cancelled previous analysis for {interview_uuid}")
              
              # Start new background analysis task
              analysis_task = asyncio.create_task(
                  self._run_background_analysis(interview_uuid, outline, history, interview_type, 
                                              interview_language, supervision_mode)
              )
              self._background_analyses[interview_uuid] = analysis_task
              logger.info(f"[BACKGROUND] Started analysis for {interview_uuid} (mode: {supervision_mode})")
              
          except Exception as e:
              logger.error(f"Error starting background analysis: {e}")
      
      async def _run_background_analysis(self, interview_uuid: str, outline: str, history: List, 
                                      interview_type: str, interview_language: str, 
                                      supervision_mode: str = "FLEXIBLE"):
          """Core logic for executing background supervision analysis"""
          try:
              logger.info(f"[BACKGROUND] Running supervisor analysis for {interview_uuid}")
              
              # Call supervisor agent for analysis
              analysis = await generate_supervisor_analysis(
                  outline, history, interview_type, interview_language, None, supervision_mode
              )
              
              # Cache analysis results
              self._cached_analyses[interview_uuid] = {
                  'analysis': analysis,
                  'timestamp': datetime.utcnow(),
                  'history_length': len(history)
              }
              
              # Update statistics
              self._analysis_count[interview_uuid] = self._analysis_count.get(interview_uuid, 0) + 1
              
              logger.info(f"[BACKGROUND] Analysis cached for {interview_uuid} (count: {self._analysis_count[interview_uuid]})")
              
          except asyncio.CancelledError:
              logger.info(f"[BACKGROUND] Analysis cancelled for {interview_uuid}")
          except Exception as e:
              logger.error(f"[BACKGROUND] Analysis failed for {interview_uuid}: {e}")
      
      async def get_cached_analysis(self, interview_uuid: str, current_history_length: int) -> Optional[str]:
          """Get cached supervision analysis results - Core performance optimization"""
          if interview_uuid not in self._cached_analyses:
              logger.info(f"[BACKGROUND] No cached analysis for {interview_uuid}")
              return None
          
          cached = self._cached_analyses[interview_uuid]
          
          # Check if cache has expired
          age = (datetime.utcnow() - cached['timestamp']).total_seconds()
          if age > self._cache_ttl:
              logger.info(f"[BACKGROUND] Cached analysis expired for {interview_uuid} (age: {age:.1f}s)")
              del self._cached_analyses[interview_uuid]
              return None
          
          # Check if conversation history has changed significantly
          history_diff = current_history_length - cached['history_length']
          if history_diff > 4:  # If conversation has progressed too much, cache may be outdated
              logger.info(f"[BACKGROUND] Cached analysis outdated for {interview_uuid} (history diff: {history_diff})")
              return None
          
          logger.info(f"[BACKGROUND] Using cached analysis for {interview_uuid} (age: {age:.1f}s)")
          return cached['analysis']
      
      def get_cost_stats(self, interview_uuid: str) -> dict:
          """Get cost statistics information - For monitoring and optimization"""
          return {
              'analysis_count': self._analysis_count.get(interview_uuid, 0),
              'estimated_tokens_saved': self._analysis_count.get(interview_uuid, 0) * (1000 - SUPERVISOR_MAX_TOKENS),
              'cache_hits': 1 if interview_uuid in self._cached_analyses else 0
          }
      
      def cleanup_interview(self, interview_uuid: str):
          """Clean up interview-related resources - Prevent memory leaks"""
          cost_stats = self.get_cost_stats(interview_uuid)
          
          # Cancel background tasks
          if interview_uuid in self._background_analyses:
              self._background_analyses[interview_uuid].cancel()
              del self._background_analyses[interview_uuid]
          
          # Clear cache
          if interview_uuid in self._cached_analyses:
              del self._cached_analyses[interview_uuid]
          
          # Clear statistics
          if interview_uuid in self._analysis_count:
              del self._analysis_count[interview_uuid]
          
          logger.info(f"[BACKGROUND] Cleaned up resources for {interview_uuid}")
          logger.info(f"[COST-STATS] {interview_uuid}: {cost_stats['analysis_count']} analyses, ~{cost_stats['estimated_tokens_saved']} tokens saved")

  # Global background supervisor manager instance
  background_supervisor = BackgroundSupervisorManager()

  # =====================================================
  # 6. Three Execution Modes - Optimization Strategies for Different Scenarios
  # =====================================================

  async def generate_ai_response_with_background(text, outline, history, interview_type, interview_language, 
                                              interview_uuid=None, supervision_mode="FLEXIBLE"):
      """Dual-Agent System - Background Mode (Fastest, Recommended)
      
      Core advantages of background mode:
      1. Asynchronous supervision analysis during user speech
      2. AI response uses cached supervision results
      3. Minimizes user waiting time
      4. Cost optimization
      """
      logger.info("=== DUAL-AGENT SYSTEM (BACKGROUND MODE) ===")
      print(f"[BACKGROUND-DUAL-AGENT] Using cost-optimized background supervision (mode: {supervision_mode})...")
      
      try:
          start_time = asyncio.get_event_loop().time()
          
          # Check if supervision analysis is needed
          use_supervisor = should_trigger_supervisor(history, interview_uuid)
          supervisor_analysis = None
          
          if use_supervisor and interview_uuid:
              # Try to get cached supervision analysis
              supervisor_analysis = await background_supervisor.get_cached_analysis(interview_uuid, len(history))
              supervision_mode_display = f"SUPERVISED-{supervision_mode}" if supervisor_analysis else f"BASIC-{supervision_mode}"
          else:
              supervision_mode_display = f"BASIC-{supervision_mode}"
          
          # Immediately start background analysis for next conversation round (Key to performance optimization)
          if interview_uuid:
              next_history = history + ([{"role": "user", "content": text}] if text else [])
              await background_supervisor.start_background_analysis(
                  interview_uuid, outline, next_history, interview_type, interview_language, supervision_mode
              )
          
          # Build system prompt (Based on whether supervision analysis is available)
          if supervisor_analysis:
              # Enhanced prompt with supervision guidance
              system_prompt = f"""You are a professional AI interviewer. Conduct the interview based on supervision analysis:

  **Supervision Analysis and Guidance**:
  {supervisor_analysis}

  **Interview Outline**: {outline}

  Please conduct the interview according to the supervision analysis guidance, ensuring high-quality conversation."""
          else:
              # Basic interview prompt
              system_prompt = f"""You are a professional AI interviewer, conducting interviews based on the following outline:

  {outline}

  Please conduct natural, in-depth interview conversations."""
          
          # Add language adaptation
          if interview_language:
              language_hints = {
                  "CHINESE": "Please conduct the interview in Chinese.",
                  "ENGLISH": "Please conduct the interview in English.",
                  "FRENCH": "Please conduct the interview in French.",
                  "NORWEGIAN": "Please conduct the interview in Norwegian.",
                  "FLEXIBLE": "Please adapt to the language used by the interviewee."
              }
              system_prompt += f"\n{language_hints.get(interview_language, '')}"
          
          # Build conversation messages
          base_messages = [*history]
          if text:
              base_messages.append({"role": "user", "content": text})
          elif not history and not text:
              # Initial message
              if interview_language == "CHINESE":
                  base_messages.append({"role": "user", "content": "Please introduce yourself and start the interview."})
              else:
                  base_messages.append({"role": "user", "content": "Please introduce yourself and start the interview."})
          
          # Parameter optimization (structured vs flexible mode)
          if supervision_mode_display.startswith("SUPERVISED-STRUCTURED") or supervision_mode_display.startswith("BASIC-STRUCTURED"):
              temperature = 0.1  # Structured interview: Low temperature, high consistency
              max_tokens = 500
          else:
              temperature = 0.7  # Flexible interview: Medium temperature, more natural
              max_tokens = 1000
          
          # Select responder agent
          resp_provider, resp_model = select_responder_model(supervision_mode)
          content = await generate_with_selected_provider(
              system_prompt,
              base_messages,
              temperature,
              max_tokens,
              resp_provider,
              resp_model
          )
          
          # Performance statistics
          total_time = asyncio.get_event_loop().time() - start_time
          cost_stats = background_supervisor.get_cost_stats(interview_uuid) if interview_uuid else {}
          
          supervisor_name = f"{SUPERVISOR_MODEL}" if claude_client else f"{AI_MODEL} (fallback)"
          logger.info(f"BACKGROUND DUAL-AGENT SUCCESS in {total_time:.2f}s ({supervision_mode_display} mode)")
          print(f"[BACKGROUND-DUAL-AGENT] Response generated in {total_time:.2f}s ({supervision_mode_display} mode)!")
          print(f"   Supervisor: {supervisor_name} ({'used' if supervisor_analysis else 'skipped'})")
          print(f"   Responder: {resp_provider}:{resp_model}")
          print(f"   Cost Stats: {cost_stats.get('analysis_count', 0)} analyses, ~{cost_stats.get('estimated_tokens_saved', 0)} tokens saved")
          
          model_info = f"BACKGROUND-DUAL-AGENT ({supervision_mode_display}): {resp_provider}:{resp_model} (supervisor: {supervisor_name}) - {total_time:.2f}s"
          
          return content, model_info
          
      except Exception as e:
          logger.error(f"Error in background dual-agent system: {e}")
          # Fallback to single-agent mode
          return await generate_single_model_response(text, outline, history, interview_type, interview_language, supervision_mode)

  async def generate_ai_response_parallel(text, outline, history, interview_type, interview_language, 
                                        supervision_mode="FLEXIBLE"):
      """Dual-Agent System - Parallel Mode (Fast)
      
      Parallel mode characteristics:
      1. Supervision analysis and basic response execute simultaneously
      2. Optimize final response through "fusion" step
      3. Balance speed and quality
      """
      logger.info("=== DUAL-AGENT AI SYSTEM (PARALLEL MODE) ===")
      logger.info(f"Running supervisor analysis and response generation in parallel (mode: {supervision_mode})")
      print(f"[PARALLEL-DUAL-AGENT] Starting concurrent processing (mode: {supervision_mode})...")
      
      try:
          # Build basic messages
          base_messages = [*history]
          if text:
              base_messages.append({"role": "user", "content": text})
          elif not history and not text:
              if interview_language == "CHINESE":
                  base_messages.append({"role": "user", "content": "Please introduce yourself and start the interview."})
              else:
                  base_messages.append({"role": "user", "content": "Please introduce yourself and start the interview."})
          
          # Parallel Task 1: Supervisor Agent Analysis
          async def supervisor_task():
              return await generate_supervisor_analysis(
                  outline, history, interview_type, interview_language, None, supervision_mode
              )
          
          # Parallel Task 2: Basic Response Agent (without supervision guidance)
          async def base_response_task():
              return await generate_single_model_response(
                  text, outline, history, interview_type, interview_language, supervision_mode
              )
          
          # Execute dual agents in parallel
          start_time = asyncio.get_event_loop().time()
          supervisor_analysis, (base_response, base_model) = await asyncio.gather(
              supervisor_task(),
              base_response_task(),
              return_exceptions=True
          )
          parallel_time = asyncio.get_event_loop().time() - start_time
          
          logger.info(f"Parallel execution completed in {parallel_time:.2f} seconds")
          print(f"[PARALLEL] Both agents completed in {parallel_time:.2f}s")
          
          # Handle exceptions
          if isinstance(supervisor_analysis, Exception):
              logger.error(f"Supervisor analysis failed: {supervisor_analysis}")
              supervisor_analysis = '{"strategic_guidance": "Continue the conversation and maintain the interview flow"}'
          
          if isinstance(base_response, Exception):
              logger.error(f"Base response generation failed: {base_response}")
              return "Sorry, AI service is temporarily unavailable. Please try again later.", AI_MODEL
          
          # Intelligent Fusion: Optimize basic response based on supervisor analysis
          logger.info("[FUSION] Applying supervisor guidance to base response...")
          
          fusion_prompt = f"""As an interview expert, please optimize the following AI response based on supervisor analysis.

  Supervisor Analysis:
  {supervisor_analysis}

  Original AI Response:
  {base_response}

  Please output the optimized response (keep it concise, only output the final response content):"""
          
          try:
              # Fusion processing
              fusion_provider, fusion_model = select_responder_model(supervision_mode)
              final_response = await generate_with_selected_provider(
                  "",
                  [{"role": "user", "content": fusion_prompt}],
                  0.3,  # Low temperature ensures fusion quality
                  500,  # Token limit
                  fusion_provider,
                  fusion_model
              )
              final_response = (final_response or '').strip()
              
              # Quality check
              if not final_response or len(final_response.strip()) < 10:
                  final_response = base_response
                  logger.warning("Fusion result too short, using base response")
              
          except Exception as fusion_error:
              logger.error(f"Fusion failed: {fusion_error}")
              final_response = base_response
          
          # Statistics information
          total_time = asyncio.get_event_loop().time() - start_time
          supervisor_name = f"{SUPERVISOR_MODEL} (Claude)" if claude_client else f"{AI_MODEL} (fallback)"
          resp_provider, resp_model = select_responder_model(supervision_mode)
          
          logger.info(f"PARALLEL DUAL-AGENT SUCCESS in {total_time:.2f}s")
          print(f"[PARALLEL-DUAL-AGENT] Complete pipeline finished in {total_time:.2f}s!")
          print(f"   Supervisor: {supervisor_name}")
          print(f"   Base Response: {resp_provider}:{resp_model}")
          print(f"   Fusion: Applied")
          
          model_info = f"PARALLEL-DUAL-AGENT: {resp_provider}:{resp_model} (supervised by {supervisor_name}) - {total_time:.2f}s"
          
          return final_response, model_info
          
      except Exception as e:
          logger.error(f"Error in parallel dual-agent system: {e}")
          return await generate_single_model_response(text, outline, history, interview_type, interview_language, supervision_mode)

  async def generate_ai_response_sequential(text, outline, history, interview_type, interview_language, 
                                          supervision_mode="FLEXIBLE"):
      """Dual-Agent System - Sequential Mode (Original implementation, most precise)
      
      Sequential mode characteristics:
      1. Execute supervisor agent analysis first
      2. Generate response based on supervision results
      3. Highest quality but relatively slower speed
      """
      logger.info("=== DUAL-AGENT AI SYSTEM (SEQUENTIAL MODE) ===")
      logger.info(f"Step 1: Supervisor analysis using {SUPERVISOR_MODEL} (mode: {supervision_mode})")
      print(f"[SEQUENTIAL-DUAL-AGENT] Starting supervisor analysis with {SUPERVISOR_MODEL} (mode: {supervision_mode})...")
      
      try:
          # Step 1: Supervisor Agent Analysis
          logger.info("Getting supervisor analysis...")
          print("[SUPERVISOR] Analyzing interview progress and providing guidance...")
          supervisor_analysis = await generate_supervisor_analysis(
              outline, history, interview_type, interview_language, None, supervision_mode
          )
          
          logger.info("Supervisor analysis completed")
          print(f"[SUPERVISOR] Analysis complete. Guidance received from {SUPERVISOR_MODEL}")
          
          # Build enhanced system prompt
          if supervision_mode == "STRUCTURED":
              base_prompt = f"""You are a professional structured interviewer who must strictly follow the outline and supervision guidance.

  **Supervision Analysis and Guidance**:
  {supervisor_analysis}

  **Interview Outline**: {outline}

  Please strictly execute the interview according to the supervision analysis guidance, ensuring adherence to the outline."""
          else:
              base_prompt = f"""You are a professional AI interviewer conducting high-quality interviews based on supervision analysis.

  **Supervision Analysis and Guidance**:
  {supervisor_analysis}

  **Interview Outline**: {outline}

  Please conduct the interview according to the professional supervision guidance."""
          
          # Add language adaptation
          if interview_language:
              language_hints = {
                  "CHINESE": "Please conduct the interview in Chinese.",
                  "ENGLISH": "Please conduct the interview in English.",
                  "FRENCH": "Please conduct the interview in French.",
                  "NORWEGIAN": "Please conduct the interview in Norwegian.",
                  "FLEXIBLE": "Please adapt to the language used by the interviewee."
              }
              base_prompt += f"\n{language_hints.get(interview_language, '')}"
          # Build messages
          base_messages = [*history]
          if text:
              base_messages.append({"role": "user", "content": text})
          if not history and not text:
              if interview_language == "CHINESE":
                  base_messages.append({"role": "user", "content": "Please introduce yourself and start the interview."})
              else:
                  base_messages.append({"role": "user", "content": "Please introduce yourself and start the interview."})
          
          # Step 2: Responder Agent Generation
          logger.info("Step 2: Generating response with supervisor guidance")
          print(f"[RESPONDER] Generating response using supervised guidance...")
          
          # Parameter optimization
          if supervision_mode == "STRUCTURED":
              temperature = 0.1
              max_tokens = 500
          else:
              temperature = 0.7
              max_tokens = 1000
          
          # Select responder agent
          resp_provider, resp_model = select_responder_model(supervision_mode)
          content = await generate_with_selected_provider(
              base_prompt,
              base_messages,
              temperature,
              max_tokens,
              resp_provider,
              resp_model
          )
          
          # Statistics information
          supervisor_name = f"{SUPERVISOR_MODEL} (Claude)" if claude_client else f"{AI_MODEL} (fallback)"
          logger.info(f"SEQUENTIAL DUAL-AGENT SUCCESS: Supervisor({supervisor_name}) + Responder({resp_model})")
          print(f"[SEQUENTIAL-DUAL-AGENT] Response generated successfully!")
          print(f"   Supervisor: {supervisor_name}")
          print(f"   Responder: {resp_provider}:{resp_model}")
          
          model_info = f"SEQUENTIAL-DUAL-AGENT: {resp_provider}:{resp_model} (supervised by {supervisor_name})"
          
          return content, model_info
          
      except Exception as e:
          logger.error(f"Error generating AI response with sequential dual-agent system: {e}")
          fallback_msg = "Sorry, AI service is temporarily unavailable. Please try again later." if interview_language == "CHINESE" else "Sorry, AI service is temporarily unavailable. Please try again later."
          return fallback_msg, AI_MODEL

  # =====================================================
  # 7. Main Entry Function - Dual-Agent System Dispatcher
  # =====================================================

  async def generate_ai_response(text, outline, history, interview_type, interview_language, 
                              interview_uuid=None, supervision_mode="FLEXIBLE"):
      """Dual-Agent AI Interview System Main Entry
      
      Select execution mode based on environment variable DUAL_MODEL_MODE:
      1. "background": Background mode (fastest, recommended)
      2. "parallel": Parallel mode (fast)
      3. "sequential": Sequential mode (most precise)
      4. "single": Single-agent mode (simplest)
      """
      if not ENABLE_DUAL_MODEL:
          return await generate_single_model_response(text, outline, history, interview_type, interview_language, supervision_mode)
      
      # Check execution mode
      DUAL_MODEL_MODE = os.getenv("DUAL_MODEL_MODE", "background").lower()
      
      if DUAL_MODEL_MODE == "background":
          # Background supervision analysis mode (fastest)
          return await generate_ai_response_with_background(
              text, outline, history, interview_type, interview_language, interview_uuid, supervision_mode
          )
      elif DUAL_MODEL_MODE == "parallel":
          # Parallel execution mode
          return await generate_ai_response_parallel(
              text, outline, history, interview_type, interview_language, supervision_mode
          )
      elif DUAL_MODEL_MODE == "sequential":
          # Sequential execution mode (original)
          return await generate_ai_response_sequential(
              text, outline, history, interview_type, interview_language, supervision_mode
          )
      else:
          # Default to background mode
          logger.warning(f"Unknown DUAL_MODEL_MODE: {DUAL_MODEL_MODE}, defaulting to background")
          return await generate_ai_response_with_background(
              text, outline, history, interview_type, interview_language, interview_uuid, supervision_mode
          )

  # =====================================================
  # 8. Single-Agent Mode - Traditional Implementation (Compatibility Support)
  # =====================================================

  async def generate_single_model_response(text, outline, history, interview_type, interview_language, 
                                        supervision_mode="FLEXIBLE"):
      """Traditional Single-Agent AI Response System (Compatibility Mode)
      
      Fallback solution when dual-agent system is unavailable
      """
      try:
          # Decide prompt strategy based on supervision mode
          if supervision_mode == "STRUCTURED":
              system_prompt = f"""You are a strict structured interviewer who must conduct interviews according to the outline step by step.

  **Interview Outline** (must strictly follow):
  {outline}

  Requirements:
  1. Strictly ask questions in outline order
  2. Only ask one question at a time
  3. Ensure coverage of all outline points
  4. Do not deviate from outline topics"""
          else:
              system_prompt = f"""You are a professional AI interviewer conducting in-depth interviews based on the following outline:

  {outline}

  Please conduct natural, flexible interview conversations and explore related topics in depth."""
          
          # Add language preference
          if interview_language:
              language_hints = {
                  "CHINESE": "Please conduct the interview in Chinese.",
                  "ENGLISH": "Please conduct the interview in English.",
                  "FRENCH": "Please conduct the interview in French.",
                  "NORWEGIAN": "Please conduct the interview in Norwegian.",
                  "FLEXIBLE": "Please adapt to the language used by the interviewee."
              }
              system_prompt += f"\n{language_hints.get(interview_language, '')}"
          
          # Build messages
          base_messages = [*history]
          if text:
              base_messages.append({"role": "user", "content": text})
          if not history and not text:
              if interview_language == "CHINESE":
                  base_messages.append({"role": "user", "content": "Please introduce yourself and start the interview."})
              else:
                  base_messages.append({"role": "user", "content": "Please introduce yourself and start the interview."})
          
          # Select model and parameters
          resp_provider, resp_model = select_responder_model(supervision_mode)
          
          # Optimize parameters
          if supervision_mode == "STRUCTURED":
              temperature = 0.1
              max_tokens = 500
          else:
              temperature = 0.7
              max_tokens = 1000
          
          # Generate response
          content = await generate_with_selected_provider(
              system_prompt,
              base_messages,
              temperature,
              max_tokens,
              resp_provider,
              resp_model
          )
          
          return content, f"{resp_provider}:{resp_model}"
          
      except Exception as e:
          logger.error(f"Error generating single-agent AI response: {e}")
          fallback_msg = "Sorry, AI service is temporarily unavailable. Please try again later." if interview_language == "CHINESE" else "Sorry, AI service is temporarily unavailable. Please try again later."
          return fallback_msg, AI_MODEL

  # =====================================================
  # 9. System Initialization and Startup Information
  # =====================================================

  def display_system_banner():
      """Display system startup banner showcasing dual-agent configuration"""
      if ENABLE_DUAL_MODEL:
          DUAL_MODEL_MODE = os.getenv("DUAL_MODEL_MODE", "background").lower()
          
          mode_info = {
              "background": ("BACKGROUND MODE", "Background supervision analysis mode (fastest)"),
              "parallel": ("PARALLEL MODE", "Parallel execution mode (fast)"),
              "sequential": ("SEQUENTIAL MODE", "Sequential execution mode (slower)")
          }
          
          mode_display, mode_desc = mode_info.get(DUAL_MODEL_MODE, ("BACKGROUND MODE", "Background supervision analysis mode (default)"))
          
          print("=" * 70)
          print("DUAL-AGENT AI SYSTEM ACTIVATED")
          print(f"Optimization Mode: {mode_display}")
          print(f"Mode Description: {mode_desc}")
          print(f"Cost Optimization: Max tokens={SUPERVISOR_MAX_TOKENS}, Frequency=1/{SUPERVISOR_CALL_FREQUENCY}")
          if claude_client:
              print(f"Supervisor Agent: {SUPERVISOR_MODEL}")
          else:
              print(f"Supervisor Agent (fallback): OpenAI {AI_MODEL}")
          print(f"Response Agent: {AI_MODEL} (OpenAI)")
          print("Tip: Use 'python configure_dual_model.py --show-config' to view settings")
          print("=" * 70)
      else:
          print("=" * 60)
          print("SINGLE-AGENT AI SYSTEM (Legacy Mode)")
          print("Tip: Use 'python configure_dual_model.py --mode background' to enable dual-agent")
          print("=" * 60)

  # =====================================================
  # 10. Usage Examples and Test Functions
  # =====================================================

  async def example_dual_agent_interview():
      """Dual-Agent system usage example"""
      print("\nDual-Agent AI Interview System Demo")
      print("=" * 50)
      
      # Example interview configuration
      outline = """
      1. Personal background introduction
      2. Work experience sharing
      3. Skills and expertise discussion
      4. Future plans outlook
      """
      
      interview_history = []
      interview_type = "SEMI_STRUCTURED"
      interview_language = "ENGLISH"
      interview_uuid = "demo-001"
      
      # Simulate interview start
      print("\nInterview starting...")
      response, model_info = await generate_ai_response(
          text=None,  # Initial message
          outline=outline,
          history=interview_history,
          interview_type=interview_type,
          interview_language=interview_language,
          interview_uuid=interview_uuid,
          supervision_mode="FLEXIBLE"
      )
      
      print(f"\nAI Interviewer: {response}")
      print(f"Model Info: {model_info}")
      
      # Simulate user response
      interview_history.extend([
          {"role": "assistant", "content": response, "timestamp": datetime.utcnow().isoformat()},
          {"role": "user", "content": "I am a software engineer with 5 years of development experience.", "timestamp": datetime.utcnow().isoformat()}
      ])
      
      # Second round of conversation
      print("\nSecond round of conversation...")
      response2, model_info2 = await generate_ai_response(
          text="I am a software engineer with 5 years of development experience.",
          outline=outline,
          history=interview_history[:-1],  # Exclude the just-added user message
          interview_type=interview_type,
          interview_language=interview_language,
          interview_uuid=interview_uuid,
          supervision_mode="FLEXIBLE"
      )
      
      print(f"\nAI Interviewer: {response2}")
      print(f"Model Info: {model_info2}")
      
      # Show cost statistics
      cost_stats = background_supervisor.get_cost_stats(interview_uuid)
      print(f"\nCost Statistics: {cost_stats}")
      
      # Clean up resources
      background_supervisor.cleanup_interview(interview_uuid)

  if __name__ == "__main__":
      # Display system banner
      display_system_banner()
      
      # Run demonstration (if this file is executed directly)
      # asyncio.run(example_dual_agent_interview())
      
      print("\nDual-Agent system code showcase completed")
      print("This system implements a dual AI architecture with Supervisor Agent + Responder Agent")
      print("Supports three execution modes: background, parallel, and sequential")
      print("Includes intelligent cost optimization and caching strategies")
      print("Can significantly improve interview quality and efficiency")

\end{lstlisting}

\end{document}
\endinput
%%
%% End of file `sample-sigconf-authordraft.tex'.
